\documentclass[11pt,reqno]{amsart}
\usepackage[T1]{fontenc}
\usepackage[utf8]{inputenc}
\usepackage{lmodern,amsmath,amssymb,amsthm,mathtools}
\usepackage{booktabs,longtable,array,enumitem,microtype}
\usepackage[margin=1.05in]{geometry}
\usepackage[dvipsnames]{xcolor}
\usepackage[colorlinks=true,linkcolor=MidnightBlue,citecolor=MidnightBlue,urlcolor=MidnightBlue]{hyperref}
\usepackage{url}
\hypersetup{pdftitle={An explicit bound of twelve for cycle-edge decompositions},
 pdfauthor={Jiacheng Shi},pdfsubject={Computer-assisted proof with Lean verification},
 pdfkeywords={graph theory, cycle decomposition, Erdos-Gallai, Lean}}
\newtheorem{theorem}{Theorem}[section]
\newtheorem{lemma}[theorem]{Lemma}
\newtheorem{proposition}[theorem]{Proposition}
\newtheorem{corollary}[theorem]{Corollary}
\theoremstyle{definition}
\newtheorem{definition}[theorem]{Definition}
\newtheorem{remark}[theorem]{Remark}
\newcommand{\ce}{\operatorname{ce}}
\newcommand{\cc}{\operatorname{cc}}
\newcommand{\rk}{\operatorname{rk}}
\newcommand{\wt}{\operatorname{wt}}
\newcommand{\cost}{\operatorname{cost}}
\newcommand{\R}{\mathbb R}
\newcommand{\N}{\mathbb N}
\newcommand{\Q}{\mathbb Q}
\newcommand{\E}{\mathbb E}
\newcommand{\Prob}{\mathbb P}
\newcommand{\ep}{\varepsilon}
\newcommand{\dd}{\mathrm d}
\newcommand{\floor}[1]{\left\lfloor #1\right\rfloor}
\newcommand{\ceil}[1]{\left\lceil #1\right\rceil}
\newcommand{\abs}[1]{\left|#1\right|}
\newcommand{\norm}[1]{\left\lVert#1\right\rVert}
\newcommand{\supp}{\operatorname{supp}}
\newcommand{\sgn}{\operatorname{sgn}}
\newcommand{\ind}{\mathbf1}
\newcommand{\mcode}[1]{\texttt{\detokenize{#1}}}
\title[Cycle--edge decompositions with coefficient twelve]{An explicit bound of twelve for cycle--edge decompositions}
\author{Jiacheng Shi}
\thanks{Copyright 2026 Jiacheng Shi. This manuscript is licensed under CC BY 4.0.}
\date{November 6, 2026}
\subjclass[2020]{05C38, 05C70; 68V15}
\keywords{cycle decomposition, Erd\H{o}s--Gallai conjecture, weighted graph decomposition, computer-assisted proof, Lean}
\newif\ifCycleVerified
\CycleVerifiedtrue
\newcommand{\CycleDependencyCount}{1\,068\,913}
\newcommand{\CycleProjectDependencyCount}{999\,202}

\begin{document}
\begin{abstract}
For a finite simple graph $G$, let $c(G)$ denote its number of connected
components, including isolated vertices. We prove that the edges of $G$ admit
a partition into at most $12(\abs{V(G)}-c(G))$ simple cycles and single edges.
The argument establishes an explicit weighted bound with cycle weight $1$
and single-edge weight $23/42$, using a finite-degree completion theorem,
three successive rates in a scale decomposition, and a renewal potential.
A finite one-vertex amplification and a minimum-weight packing convert the
weighted estimate into the stated component-rank bound. Rational certificates
supply the quantitative interfaces, including all stopping phases and the
infinite tails. A complete Lean formalization of the final graph-theoretic statement
accompanies the paper; the retained project source closure was rebuilt
and its proof dependencies audited.

\end{abstract}
\maketitle
\tableofcontents
\section{Introduction}\label{sec:introduction}
All graphs in this paper are finite, simple, and undirected. A
\emph{cycle--edge decomposition} is a partition of the original edge set
into edge sets of simple cycles and singleton edges. Let $\ce(G)$ be the
minimum number of parts in such a partition. We write
\[
 n=\abs{V(G)},\qquad m=\abs{E(G)},\qquad r(G)=n-c(G),
\]
where $c(G)$ counts all connected components, including isolated vertices.
The empty graph on no vertices has $c(G)=0$.

\begin{theorem}\label{thm:main}
Every finite simple graph $G$ satisfies
\[
                 \ce(G)\le 12\bigl(\abs{V(G)}-c(G)\bigr).
\]
\end{theorem}

The problem of obtaining a linear bound for $\ce(G)$ goes back to Erd\H{o}s
and Gallai; it appears in the discussion of Erd\H{o}s, Goodman and
P\'osa~\cite{EGP1966}. Conlon, Fox and Sudakov~\cite{CFS2014} proved an
$O(n\log\log n)$ bound, and Buci\'c and Montgomery~\cite{BM2024} subsequently
obtained $O(n\log^*n)$. Linear bounds were presented in 2026 by
OpenAI~\cite{OAI2026}, by Kim~\cite{Kim2026}, and in a separate manuscript
and formalization of Coffey~\cite{Coffey2026}. The present paper addresses an
explicit coefficient. Its large-scale construction follows the
quotient-and-lifting approach of~\cite{OAI2026}; the result is not a new proof
of the mere existence of some linear constant.

Our proof first assigns weight $1$ to a cycle and weight
$\sigma=23/42$ to a singleton edge. With $C=485/42$, the central estimate is
\begin{equation}\label{eq:weighted-order}
 \exists\mathcal P\quad \wt_\sigma(\mathcal P)\le Cn,
\end{equation}
where $\mathcal P$ is an actual cycle packing and all uncovered original
edges are counted as singletons. A component-rank normalization and the
identity
\[
                  C+1-\sigma=12
\]
then give Theorem~\ref{thm:main}. Keeping the singleton weight throughout
the proof is essential: replacing weighted completion costs by unweighted
ones at intermediate stages loses the final coefficient.

The quantitative proof has three parts. First, a mixed cycle-packing estimate
with constants $43/10$ and $3$ is combined with a short-circumference base
and finite depth-first-search chains. Second, an indexed scale decomposition
uses three final degree-reduction rates and preserves the exact assignment of
every original edge. Third, a positive renewal potential pays the accumulated
losses. Forty stopping-phase intervals determine different prefix thresholds;
a positive second-order recurrence controls every possible number of layers.

The proof is computer-assisted. Two local interfaces require substantial
finite certificates: the short-circumference base and the quantitative
mixed-cover estimate. Their mathematical specifications and certification
procedures are given in Sections~\ref{sec:finite}, \ref{app:terminal}, and
\ref{app:cover}. The subsequent scale and cost calculations are displayed
explicitly, with the rational data and the complete formal proofs supplied
in the companion repository. Appendix~\ref{app:correspondence} records the
correspondence between the principal statements here and their Lean
counterparts. The coefficient throughout this paper is exactly $12$; no
parameters from later numerical experiments are used.

\section{Weighted packings and component-rank normalization}\label{sec:weighted}
A \emph{cycle packing} $\mathcal P$ is a family of pairwise edge-disjoint
simple cycles in $G$. Let $p(\mathcal P)$ be its number of cycles and let
$q(\mathcal P)$ be the number of edges not used by the cycles. Define
\begin{equation}\label{eq:weights}
 \wt_\sigma(\mathcal P)=p(\mathcal P)+\frac{23}{42}q(\mathcal P),
 \qquad
 z(\mathcal P)=42p(\mathcal P)+23q(\mathcal P).
\end{equation}
The cycles together with the $q(\mathcal P)$ remaining singleton edges form
an exact decomposition of $E(G)$. We use the empty packing when $G$ has no
cycles. Since the possible integer costs form a nonempty subset of $\N$,
a minimum-cost packing exists.

\begin{lemma}\label{lem:forest}
For a minimum-cost packing $\mathcal P$, the graph formed by its uncovered
edges is a forest. Consequently $q(\mathcal P)\le r(G)$.
\end{lemma}
\begin{proof}
If the uncovered graph contains a simple cycle of length $\ell\ge3$, adding
that cycle to the packing changes its integer cost by $42-23\ell<0$,
contrary to minimality. The uncovered graph $F$ retains all $n$ vertices,
so $q(\mathcal P)=\abs{E(F)}=n-c(F)\le n-c(G)$, since deletion of edges
cannot merge components.
\end{proof}

The following finite amplification is the normalization used by the
formal proof. In particular, no limiting argument is needed.
\begin{lemma}\label{lem:amplification}
Suppose every finite simple graph $H$ has a packing of integer cost at most
$485\abs{V(H)}$. Then every nonempty finite simple graph $G$ has a packing
of integer cost at most $485(\abs{V(G)}-1)$.
\end{lemma}
\begin{proof}
Choose a vertex $v\in V(G)$ and form a graph $W$ from $486$ disjoint copies
of $G$ by identifying their copies of $v$. The resulting graph is simple and
has $486(n-1)+1$ vertices. A simple cycle of $W$ lies entirely in a single
copy: to visit two copies it would have to pass through the common vertex
twice before closing. Every singleton edge also belongs to one copy.
Thus restriction of any packing of $W$ gives a decomposition of each copy.
If $z$ is the minimum integer cost for $G$, every packing of $W$ has cost
at least $486z$.

The assumed order bound gives
\[
 486z\le485\bigl(486(n-1)+1\bigr).
\]
If $z>485(n-1)$, integrality implies $z\ge485(n-1)+1$; multiplying by
$486$ contradicts the preceding inequality. Therefore
$z\le485(n-1)$, as required.
\end{proof}

\begin{proposition}\label{prop:weighted-reduction}
The weighted order bound~\eqref{eq:weighted-order} implies
Theorem~\ref{thm:main}.
\end{proposition}
\begin{proof}
Multiplying~\eqref{eq:weighted-order} by $42$ gives the hypothesis of
Lemma~\ref{lem:amplification}. Apply that lemma to each connected component
of $G$ and assemble the resulting packings. This gives a packing of integer
cost at most $485r(G)$. Choose a minimum-cost packing $\mathcal P$ of $G$;
then $z(\mathcal P)\le485r(G)$ and Lemma~\ref{lem:forest} gives
$q(\mathcal P)\le r(G)$. Hence
\[
 42\bigl(p(\mathcal P)+q(\mathcal P)\bigr)
   =z(\mathcal P)+19q(\mathcal P)
   \le504r(G).
\]
Dividing by $42$ proves the assertion. Isolated components contribute zero
to the rank and require no parts. If $n=0$, the empty decomposition suffices.
\end{proof}

We shall prove a stronger form of~\eqref{eq:weighted-order}, with an
order-dependent factor $\phi(n)\le1$. Every recursive call in that proof
is on a graph with strictly fewer vertices; induced subgraphs, quotient
graphs, and padded graphs are transported along explicit vertex bijections
when necessary.

\section{Classical inputs and their formal proofs}\label{sec:classical-inputs}
We record two inputs whose formal implementations are supplied in the
artifact. The first uses the rotation-and-deletion argument attributed to
Woodall~\cite{Woodall1972}; see also~\cite[Theorem 2.7.1]{Gunderson2022}.
The second is a lower-tail estimate proved from finite product weights,
the FKG inequality~\cite{FKG1971}, and a Laplace-transform argument.
These are implementations of classical methods, not claims of new
extremal or correlation inequalities.

\begin{lemma}[Erd\H{o}s--Gallai cycle bound]\label{lem:eg-input}
Let $G$ be a finite simple graph on $n$ vertices and let $k\ge2$ be an
integer. If every simple cycle has length at most $k$, then
\[
                         2\abs{E(G)}\le k(n-1)
\]
when $n\ge1$. For $n=0$ both sides of the corresponding natural-number
statement are zero, with the predecessor interpreted as $n\mathbin{\dot-}1$.
\end{lemma}
\begin{proof}
We use strong induction on $n$. The edgeless case is immediate. Otherwise
choose a longest path $P=v_0\cdots v_\ell$, maximizing $d=\deg(v_0)$ among
the starting degrees of all longest paths. Then $\ell\ge1$. All neighbors
of the first vertex of any longest path lie on that path, since an outside
neighbor would extend it.

Define
\[
 W=\{v_i:0\le i<\ell,\ v_0v_{i+1}\in E(G)\},\qquad
 X=\{v_j:0\le j<\min\{k,\ell+1\}\}.
\]
The predecessor operation is a bijection from $N(v_0)$ to $W$, so
$\abs W=d\ge1$, and $v_\ell\notin W$. For each $v_i\in W$, the rotation
\[
             v_i v_{i-1}\cdots v_0v_{i+1}\cdots v_\ell
\]
is another longest path with the same vertex set. Consequently
$\deg(v_i)\le d$ and $N(v_i)\subseteq V(P)$.

A chord $v_0v_j$ with $j\ge2$ closes a prefix of $P$ to a simple cycle
of length $j+1$, whence $j<k$. The case $j=1$ also has $j<k$.
In particular, $v_i\in W$ implies $i+1<k$, and $W\subseteq X$.
If $v_i\in W$ has a neighbor $v_j$ with $j\ge k$, then $i+1<j$ and
\[
                v_0\cdots v_i v_j v_{j-1}\cdots v_{i+1}v_0
\]
is a simple cycle of length $j+1>k$. Thus every neighbor of $W$ lies in $X$.

Let $a$ count the edges inside $W$ and let $b$ count the edges between
$W$ and $X\setminus W$. The degree bounds and simplicity give
\[
             2a+b\le d^2,\qquad b\le d(\abs X-d).
\]
The $a+b$ edges removed with $W$ therefore satisfy
\[
                     2(a+b)\le d\abs X\le dk.
\]
The remaining graph has $n-d$ vertices, where $1\le n-d<n$, and retains
the same bound on cycle lengths. Applying induction and adding the
deletion cost yields
\[
 2\abs{E(G)}\le dk+k(n-d-1)=k(n-1).
\]
No connectivity hypothesis is used.
\end{proof}

In Lean, the rotations and the degree-maximizing choice are implemented in
\begin{center}\small\nolinkurl{Cycle12/Independent/PathRotation.lean}.\end{center}
The set $W$ and its neighborhood bounds are implemented in
\begin{center}\small\nolinkurl{Cycle12/Independent/WoodallSet.lean}.\end{center}
The counting step is implemented by splitting degree sums across $W$ and
its complement, so internal edges are counted exactly twice in the degree
sum and exactly once in the deletion. The final strong induction is
\mcode{Cycle12.Independent.erdosGallai_cycle}. Passing to a longest cycle
gives the precise long-cycle interface used in the finite completion
argument; its denominator is $n-1$.

\begin{lemma}[Lower tail for a Bernoulli vertex sample]\label{lem:sampling-input}
Let $G$ have $m$ edges, let $D\ge0$, and suppose that
$\deg(u)+\deg(v)\le D+2$ for every edge $uv$.
Independently select each vertex with probability $p\in[0,1]$, and let $X$
count the edges with both endpoints selected. Put
\[
                       \mu=mp^2,\qquad v=mp^2+mDp^3.
\]
If $\mu>0$ and $0\le t\le\mu$, then
\[
                    \Prob(X\le\mu-t)\le
                    \exp\!\left(-\frac{t^2}{2v}\right).
\]
If $\mu=0$, the only required value is $t=0$, for which the bound is the
elementary probability bound by $1$.
\end{lemma}
\begin{proof}
For a subset $S$ of the vertex set, its probability mass is
$w(S)=p^{\abs S}(1-p)^{n-\abs S}$. The identity
\[
                     w(A)w(B)=w(A\cap B)w(A\cup B)
\]
follows by adding the corresponding cardinalities and complement
cardinalities. FKG gives positive association of nonnegative increasing
functions. Shifting an increasing function by its value at the empty set
extends this to real-valued increasing functions. Applying positive
association to $f$ and $M-g$, where $g$ is decreasing and $M\ge\max g$,
gives
\begin{equation}\label{eq:opposite-association}
                 \E(fg)\le (\E f)(\E g)
\end{equation}
for nonnegative increasing $f$ and nonnegative decreasing $g$.
Independence of observables using disjoint coordinates follows by
induction on the finite ground set, conditioning on its last coordinate.

Write $I_e$ for the indicator that edge $e$ survives. For a fixed edge
$e$, split
\[
                         X=I_e+Y_e+Z_e,
\]
where $Y_e$ sums the indicators of the other edges sharing an endpoint
with $e$, and $Z_e$ sums all remaining indicators. There are at most $D$
terms in $Y_e$. Coordinate independence gives, for $s\ge0$,
\[
                  \E(I_e e^{-sZ_e})=p^2\E(e^{-sZ_e}).
\]
For each edge $f$ contributing to $Y_e$, the increasing function $I_eI_f$
has expectation $p^3$; the formal estimate only needs the upper bound
$p^3$. Applying \eqref{eq:opposite-association} and summing gives
\[
                  \E(I_eY_e e^{-sZ_e})\le Dp^3\E(e^{-sZ_e}).
\]
The inequality $e^{-sY_e}\ge1-sY_e$, applied on the event $I_e=1$, now gives
\[
 \E(I_e e^{-sX})\ge e^{-s}(p^2-sDp^3)\E(e^{-sZ_e}).
\]
If $0\le s\le(1+Dp)^{-1}$, the coefficient on the right is nonnegative.
Since $Z_e\le X$, putting $L(s)=\E(e^{-sX})$ yields
\[
 \E(I_e e^{-sX})\ge e^{-s}(p^2-sDp^3)L(s)
                    \ge p^2\bigl(1-s(1+Dp)\bigr)L(s).
\]
For the second inequality, use $e^{-s}\ge1-s$ and then discard the
nonnegative term $s^2Dp^3$. Summation over $e$ gives
\[
                        \E(Xe^{-sX})\ge(\mu-vs)L(s).
\]
Finite summation permits differentiation, with
$L'(s)=-\E(Xe^{-sX})$. Hence
\[
                 s\longmapsto e^{\mu s-vs^2/2}L(s)
\]
is nonincreasing on $[0,\mu/v]$. Its value at $0$ is $1$, so
$L(s)\le e^{-\mu s+vs^2/2}$ there. Finally, the exponential Markov
inequality applied with $s=t/v$ proves the claim. The case $t=0$ is
handled directly by the probability bound by $1$.
\end{proof}

The finite-law calculation is formalized first for a family of distinct
two-element sets, with explicit dependency neighborhoods. The graph
interface instantiates these sets with the endpoints of the original
edges and uses the degree bound in the line graph. The statements
{\small\nolinkurl{Cycle13.vertex_sample_edge_lower_tail}} and
{\small\nolinkurl{Cycle13.finite_graph_janson_restricted}} retain the exact constants
and the finite-sample conventions required later. The proof invokes
Mathlib's FKG theorem and does not assume a separate Janson theorem.

\section{Parameters and the renewal potential}\label{sec:potential}
All constants in this section are fixed rational numbers. Their exact values
are listed in Appendix~\ref{app:parameters}; the symbols agree with the
formal proof, except that $c_0$ denotes its constant \mcode{c}, $B_0$ its
constant \mcode{batch}, and $\chi$ its constant \mcode{overlap}.
In particular
\[
 C=\frac{485}{42},\quad \sigma=\frac{23}{42},\quad
 \kappa=\frac{999999}{1000000},\quad
 \phi_0=\frac{46220729}{50000000},\quad
 B_{\rm main}=\frac{158387}{1000}.
\]
The parameters satisfy
\begin{equation}\label{eq:parameter-order}
 0<s<\theta_2\le\theta_1\le\theta_0<1,\quad 0<b<1,\quad
 \rho=1+\delta+\ep,\quad q_0=\theta_1\theta_2,\quad q_0w_A=w_0.
\end{equation}
Put
\[
 h=\frac{3s-\rho-2\alpha}{2}>0,
 \qquad c(w)=c_0\,2^{h(w-w_A)}.
\]
This growing team scale will be used in Section~\ref{sec:routing}.
All logarithms in exponents and scale coordinates are to base $2$;
$\ln$ denotes the natural logarithm.

For $z>0$, define
\begin{equation}\label{eq:scale-series}
 E(z)=\sum_{j\ge0}2^{-z\theta_2^{-j}}.
\end{equation}
The series converges. Indeed
$\theta_2^{-j}\ge1+j(\theta_2^{-1}-1)$, which bounds it above by a
geometric series. The losses at a scale $w\ge w_A$ are
\begin{align}
 F(w)&=\left(\frac{1001}{1000}+(\rho-s)q_0w\right)
                   2^{-(\alpha-sq_0)w},\label{eq:F}\\
 O(w)&=\chi M E(\delta\theta_2^2w),&
 N(w)&=M^2F(w),\label{eq:ON}\\
 J(w)&=\frac M2 E(\ep\theta_2^2w),&
 U(w)&=M2^{-(s-\lambda)w}+2M^2 2^{-(s+b-\lambda-\rho)w},\label{eq:JU}\\
 V(w)&=\frac{\gamma M}{2c(w)}F(w).\label{eq:V}
\end{align}
The total positive loss used for coefficient $12$ is
\begin{equation}\label{eq:H}
 H(w)=O(w)+\kappa\bigl(N(w)+U(w)\bigr)
                   +\frac{B_0O(w)+J(w)+V(w)}C.
\end{equation}
No cancellation between separately bounded positive losses is used in this
potential.

\begin{lemma}\label{lem:loss-decay}
Let $r_* =\ep\theta_2^2$. Each loss in~\eqref{eq:ON}--\eqref{eq:V} is
nonnegative and nonincreasing on $[w_A,\infty)$. Moreover
\begin{equation}\label{eq:H-decay}
 H(y)\le 2^{-r_*(y-x)}H(x)\qquad(w_A\le x\le y).
\end{equation}
\end{lemma}
\begin{proof}
The rational parameter inequalities imply
\[
 0<r_*\le\min\{\delta\theta_2^2,\alpha-\theta_0q_0,
 s-\lambda,s+b-\lambda-\rho,\alpha-\theta_0q_0+h\}.
\]
Every summand of $E(aw)$ decays at least as fast as $2^{-aw}$.
For $F$, put $A=1001/1000$ and $B=(\rho-s)q_0$. The certified entropy
slope inequality at $w_A$ is
\[
 B\le (\ln2)(\theta_0-s)q_0(A+Bw_A).
\]
Since $B\ge0$, its right-hand side increases with $w$. Differentiating
$\ln F(w)$ therefore gives
\[
 \frac{F'(w)}{F(w)}=\frac B{A+Bw}-(\ln2)(\alpha-sq_0)
       \le-(\ln2)(\alpha-\theta_0q_0).
\]
Integrating gives the asserted exponential comparison for $F$ and $N$.
The additional factor $2^{-h(w-w_A)}$ gives the displayed rate for $V$.
The two summands of $U$ and the series defining $O,J$ have the other rates
listed above. Positive linear combinations preserve the comparison with
the common rate $r_*$, proving~\eqref{eq:H-decay}.
\end{proof}

Define, for $x\ge bw_A$,
\begin{equation}\label{eq:renewal-loss}
 L(x)=\sum_{j\ge1}H(x/b^j).
\end{equation}
Lemma~\ref{lem:loss-decay} implies convergence, by the same geometric
comparison as for~\eqref{eq:scale-series}. Shifting the summation index gives
\begin{equation}\label{eq:renewal-identity}
                         L(bw)-L(w)=H(w)\qquad(w\ge w_A).
\end{equation}
The anchor certificate proves
\begin{equation}\label{eq:anchor-floor}
                         L(bw_A)\le1-\phi_0.
\end{equation}
Set
\begin{equation}\label{eq:Phi}
 \Phi(x)=\begin{cases}\phi_0,&x<bw_A,\\1-L(x),&x\ge bw_A,\end{cases}
 \qquad
 \phi(n)=\begin{cases}\phi_0,&n=0,\\\Phi(\log_2n),&n>0.\end{cases}
\end{equation}
The separate definition at zero matches the value used in Lean.

\begin{lemma}\label{lem:potential}
The functions $\Phi$ and $\phi$ are nondecreasing, and
$0<\phi_0\le\Phi(x)\le1$. If $w\ge w_A$ and $n\ge2^w$, then
\begin{equation}\label{eq:potential-payment}
                   H(w)\le\phi(n)-\Phi(bw).
\end{equation}
\end{lemma}
\begin{proof}
The summands in~\eqref{eq:renewal-loss} are nonnegative and nonincreasing,
so $L$ is nonnegative and nonincreasing. At $bw_A$, inequality
\eqref{eq:anchor-floor} prevents a downward jump in~\eqref{eq:Phi}.
This proves the bounds and monotonicity. Equation~\eqref{eq:renewal-identity}
gives $\Phi(w)-\Phi(bw)=H(w)$, and
$\log_2n\ge w$ gives~\eqref{eq:potential-payment}.
\end{proof}

For later use define the batch and group fees
\begin{equation}\label{eq:batch-group-fees}
 B(w)=\frac{2^{(\rho-s-\alpha)w}}{4c(w)}
                +\frac12\,2^{(\rho-s-\lambda)w},
 \qquad g(w)=\frac1{2c(w)}.
\end{equation}
They are nonnegative and nonincreasing. The endpoint certificates establish,
for every $w\ge w_A$,
\begin{equation}\label{eq:fee-endpoints}
          B(w)\le B_0,\qquad g(w)\le C\kappa\phi_0
                                  \le C\kappa\Phi(bw).
\end{equation}
All numerical assertions in this section are rational inequalities or
certified inequalities involving powers of $2$. Section~\ref{sec:formal}
and Appendix~\ref{app:parameters} specify how they are checked.

\section{Finite-degree completion}\label{sec:finite}
For real $a$ and positive $t$, write $\mathcal A(a,t)$ for the assertion
that every finite simple graph with $n$ vertices and $m$ edges has a packing
of weight at most $an+m/t$. This section supplies the finite-degree cases
of the main induction and the completion of the stopped remainder.

\subsection{The short-circumference starting point}
\begin{lemma}[Short-circumference certificate]\label{lem:six}
If $G$ has no simple cycle of length at least $7$, then it has a packing
$\mathcal P$ satisfying
\[
                  z(\mathcal P)\le19r(G)+6m.
\]
\end{lemma}
The proof is a structural reduction followed by a finite certificate check;
it is given in Appendix~\ref{app:terminal}. Both the reduction to the finite
catalogue and the interpretation of its witnesses as actual original-edge
packings are part of the formal proof.

\begin{lemma}\label{lem:harmonic}
For every integer $u\ge7$,
\begin{equation}\label{eq:harmonic}
 \mathcal A\left(\frac{19}{42}+\frac12\sum_{j=8}^u\frac1j,u\right)
 \quad\hbox{holds}.
\end{equation}
\end{lemma}
\begin{proof}
Greedily delete simple cycles of length at least $7$. Their number is at
most one seventh of the number of deleted edges. On the remaining graph,
Lemma~\ref{lem:six}, divided by $42$, gives weight at most
$19n/42+m'/7$. Adding the deleted cycles proves $\mathcal A(19/42,7)$.
For the successor step from $t$ to $t+1$, first remove cycles of length at
least $t+1$. On the remaining graph, the Erd\H{o}s--Gallai long-cycle edge
bound~\cite{EG1959} gives $2m'\le tn$. Consequently
\[
 \frac{m'}t-\frac{m'}{t+1}
          =\frac{m'}{t(t+1)}\le\frac{n}{2(t+1)}.
\]
Apply $\mathcal A(a,t)$ to this remainder and pay the deleted cycles by
their lengths at denominator $t+1$. This proves
$\mathcal A(a+1/(2(t+1)),t+1)$.
Iterating that step from $7$ to $u$ gives~\eqref{eq:harmonic}.
\end{proof}

\subsection{The mixed cover}
Put $B_*=43/10$, and define
\[
 D(n,m)=B_*nm^{1/3}+3n^{3/2},\qquad
 \psi(v)=\begin{cases}(53/200)v,&0\le v\le1,\\
              1/2-(47/200)v^{-1/3},&v\ge1.
              \end{cases}
\]
The two formulas for $\psi$ agree at $1$. The function is nonnegative,
nondecreasing and concave, and is at most $1/2$.

\begin{lemma}[Mixed-cover certificate]\label{lem:mixed}
Every graph of positive order has each of the following packings:
\begin{align}
 p&\le n\psi\left(\frac{m}{2B_*^3n}\right),
       &q&\le D(n,m)+2n;\label{eq:fine-cover}\\
 p&\le n/2,&q&\le D(n,m)+n.\label{eq:coarse-cover}
\end{align}
The packings in the two lines need not be the same.
\end{lemma}
\begin{proof}
Use strong induction on the order. If $m\le D(n,m)$, the empty packing
works. The greedy range is handled by the long-cycle estimate. A sparse
cut is paid from the difference between the parent budget $D(n,m)$ and
the two child budgets; cycles in the two induced sides are assembled and
the cut edges are left as singletons. The perspective of the concave
function $\psi$ pays the cycle counts in~\eqref{eq:fine-cover}.
A sufficiently low-degree vertex is handled by deletion and the corresponding
one-vertex budget inequality. The remaining graphs meet the degree and cut
hypotheses of the routed-sampling construction. Its certificate covers
\emph{all} the remaining continuous parameters, including the unbounded
scale range. Appendix~\ref{app:cover} gives the construction, the precise
coverage interface, and the finite checks closing this induction.
\end{proof}

\subsection{A certified depth-first-search step}
Consider a depth-first-search block with $s$ new vertices and at most $T$
boundary vertices, and let $m'$ denote its assigned edges. The boundary is
padded to $T$ vertices when needed. Choose the local packings first, and
then apply $\mathcal A(a,t)$ once to the graph of all uncovered original
edges, retaining the original vertex set. Its intercept is therefore $an$.
For a regular block, the fine packing contributes the following additional
fee, after subtracting the block's assigned-edge budget $m'/T$:
\begin{equation}\label{eq:regular-fee}
 (s+T)\psi\left(\frac{m'}{2B_*^3(s+T)}\right)
       +\frac{D(s+T,m')+2(s+T)}t-\frac{m'}T.
\end{equation}
The alternative is to leave the block's edges for that common suffix
estimate; its additional fee is $m'/t-m'/T$. An entire component with no boundary instead uses
\eqref{eq:coarse-cover}, giving fee
\begin{equation}\label{eq:component-fee}
                  s/2+\frac{D(s,m')+s}t-\frac{m'}T.
\end{equation}
These are different branches of the DFS decomposition, not interchangeable
versions of the same packing.

For $g>0$ let $\Pi_p(g)=(1+g)^pg^{1-p}$. Write
\begin{align}
 \mathcal K(T,t,d,g)
  ={}&(1+g)\left[\psi\left(\frac{d}{4B_*^3}\right)
                              +\frac2t-\frac d{2T}\right]\notag\\
     &+\frac{B_*}{t}\left(\frac{dT}{2}\right)^{1/3}\Pi_{4/3}(g)
                +\frac{3\sqrt T}{t}\Pi_{3/2}(g).\label{eq:core-fee}
\end{align}
The substitution $d=2m'/(s+T)$ and $g=T/s$ makes
\eqref{eq:regular-fee} equal to $s\mathcal K(T,t,d,g)$.
The functions $\Pi_p$ are convex for $p\ge1$: they are perspectives of
$x\mapsto x^p$. Thus, for fixed $d$, the maximum of~\eqref{eq:core-fee}
on a closed $g$-interval occurs at an endpoint, even when its linear
coefficient is negative.

\begin{lemma}[Finite DFS step]\label{lem:dfs-step}
Let $T,t$ be positive integers with $T\ge2^{20}$, and let $\beta_0,K>0$.
Suppose
\begin{align*}
 &(1+1/\beta_0)/2\le K,\\
 &\mathcal K(T,t,d,g)\le K\quad(d>t),\quad
       g\in\{(2\beta_0+2^{-20})^{-1},\beta_0^{-1}\},\\
 &s/2+\bigl(D(s,m')+s\bigr)/t-m'/T\le Ks
                 \quad(0<s\le\beta_0T+1,\ m'\ge0).
\end{align*}
Then $\mathcal A(a,t)$ implies $\mathcal A(a+K,T)$.
\end{lemma}
\begin{proof}
First remove simple cycles of length at least $T$. If $h$ edges are
removed, these cycles cost at most $h/T$. The remaining graph, on the same
vertex set, has $m-h$ edges and no cycle of length at least $T$; this is the
circumference hypothesis required by the DFS block construction.
The rounded DFS construction assigns each remaining edge to exactly one block.
Its regular blocks have $T/s$ between the two displayed endpoints; the
rounding error is bounded by $1/T\le2^{-20}$. For $d\le t$, the direct
suffix branch is paid by $(1+1/\beta_0)/2\le K$. For $d>t$, convexity
and the two endpoint hypotheses pay~\eqref{eq:regular-fee}. Boundary-free
components have at most $\beta_0T+1$ vertices and are paid by the last
hypothesis. Sum over the disjoint sets of new vertices. The assigned edge
counts sum to $m-h$. Apply the suffix packing once to the actual residual
graph, so its intercept is charged once per original vertex. This gives
weight at most $(a+K)n+(m-h)/T$ on the remainder. Restoring the initially
removed cycles adds at most $h/T$, yielding $(a+K)n+m/T$.
The original-edge assembly is proved by
\mcode{Cycle13.weightedAffine_dfs_step}; the scalar hypotheses are supplied in
\mcode{Cycle12.FiniteComponentPayment}.
\end{proof}

The endpoint hypotheses in Lemma~\ref{lem:dfs-step} are verified on whole
real density intervals. In the high-density branch, $t\ge4B_*^3=79507/250$.
After substitution, the nonlinear density expression has the form
$A x^{1/3}-B x^{-1/3}-x/2$; a certified tangent and derivative sign pay its
entire domain. Rational lower and upper bounds for $A,B$ and $\sqrt{\cdot}$
are established by cubing or squaring nonnegative rational endpoints.
The component branch has $B=0$ and also includes zero density.
No finite sampling of possible graph densities is used.

There are $83$ finite chains. A chain begins with~\eqref{eq:harmonic}, uses
a certified upper bound on its harmonic intercept, and applies
Lemma~\ref{lem:dfs-step} a finite number of times. If its final affine bound
is $\mathcal A(a,T)$ and its degree cap is $d_0$, then
\begin{equation}\label{eq:finite-completion}
       2m\le d_0n\quad\Longrightarrow\quad
       \wt_\sigma(\mathcal P)\le\left(a+\frac{d_0}{2T}\right)n.
\end{equation}
All $T$ and $d_0$ in the formal certificates are integers; bounds such as
$2^x\le d_0$ are certified before~\eqref{eq:finite-completion} is applied.
The chains and their exact fees are tabulated in the companion data.

\begin{proposition}[Finite potential base]\label{prop:finite-base}
If $2m\le2^{B_{\rm main}}n$, then $G$ has a packing of weight at most
$C\phi(n)n$.
\end{proposition}
\begin{proof}
The $45$ base intervals cover the required order range, with their shared
endpoints included. On each interval choose its certified finite chain.
Use the simple-graph estimate $2m\le n^2$ together with the assumed cap
to establish the chain's integer degree cap. The resulting constant
in~\eqref{eq:finite-completion} is at most the lower endpoint value of
$C\phi(n)$. Monotonicity from Lemma~\ref{lem:potential} pays the entire
interval. Beyond the last interval the main finite chain and the same
monotonicity apply. The zero-order case has no edges. The endpoint payments
and the interval-coverage proof are
\mcode{Cycle12.FinitePotentialPayments} and
\mcode{Cycle12.FinitePotentialBase12}.
\end{proof}

\section{An exact decomposition into scales}\label{sec:layers}
The large-order proof uses the scale decomposition of~\cite{OAI2026}, with
explicit parameters, a parity correction, and three final reduction rates.
We specify which edges and which vertex multiplicities are counted.

A supported piece consists of a graph on the original vertex set together
with a vertex support containing both endpoints of every one of its edges.
Its mass is the size of that support. Different pieces may share vertices;
their edge sets must be disjoint. The mass of a family is the sum of its
support sizes, including multiplicity.

\begin{lemma}[One scale]\label{lem:scale}
At logarithmic scale $w$, set
\[
 D=2^w,\quad \ell=2^{(1+\ep)w},\quad L_0=2^{\rho w},\quad h_0=2^{sw}.
\]
Under the parameter inequalities at the applicable scale, a piece with
average degree at most $D$ can be partitioned into long cycles, even
expanding cores, and a new family of supported remainder pieces. The
following bounds hold per unit parent mass:
\begin{enumerate}[label=(\roman*),leftmargin=*]
\item the number of long cycles is at most $D/(2\ell)$;
\item the sum of box orders is at most
      $1+2\ell/(L_0-\ell-2)$, and each box has order at most $L_0$;
\item cores in a box have disjoint vertex supports, orders greater than
      $h_0$, and cut expansion at least $(1023/1024)h_0$; every core degree
      is even;
\item the next average-degree bound is at most
      $h_0(1001/1000+\log_2(L_0/h_0))$.
\end{enumerate}
\end{lemma}
\begin{proof}
First remove actual edge-disjoint cycles longer than the prescribed
threshold; their lengths pay (i). On the cycle-free remainder, a DFS
partition with integer parameters $\floor\ell$ and
$\floor{(L_0-\floor\ell)/2}$ gives the boxes in (ii). In particular, the
denominator includes both rounding units.
Within each box, repeatedly cut a side of insufficient edge expansion.
The entropy potential pays the cut edges and gives the logarithmic term
in (iv). The surviving induced cores have disjoint supports. For each core,
choose a forest whose odd-degree vertices are exactly the odd-degree
vertices of the core, and move that forest into the remainder. Subtraction
makes the core even. The cut bound for the forest and the certified scale
inequalities leave the expansion in (iii); its at most one edge per vertex
is absorbed by the $1001/1000$ term in (iv).
The remainder children are defined \emph{after} this move. Each operation
is a partition of original edges, so composition preserves the partition.
The conditions $\ell+2<L_0$, $h_0>4096$, $h_0\le L_0$, and
$2048L_0<h_0^2$ are checked on the complete scale ranges, not assumed from
asymptotic notation.
\end{proof}

Let $X=\log_2n\ge w_A$. Choose $k\ge0$ so that
\begin{equation}\label{eq:phase}
           w_A\le W:=\theta_2^kX<w_A/\theta_2,
 \qquad m_0=k+3.
\end{equation}
(The symbol $m_0$ counts stages, not edges.) Define
\begin{equation}\label{eq:schedule}
 w_i=\begin{cases}
       \theta_2^iX,&0\le i\le k+1,\\
       \theta_1\theta_2^{k+1}X,&i=k+2,\\
       \theta_0\theta_1\theta_2^{k+1}X,&i=k+3.
      \end{cases}
\end{equation}
Apply Lemma~\ref{lem:scale} at stages $0,\ldots,k+2$. The entropy
inequalities for the three transitions give average degree at most
$2^{w_{i+1}}$ in every child piece. The last processed scale is at least
$q_0w_A=w_0$. Let $S_i$ be the total core mass at stage $i$.

For each prefix, the long cycles removed so far, the cores extracted so far,
and the \emph{whole} current frontier form an exact partition of $E(G)$.
This prefix identity permits both operations used later: replacing one
selected core layer together with its descendant frontier, or stopping
before the last two stages. It would not suffice to retain only the sparse
edges at the end of all stages.

The overlap factor in (ii) is at most $1+\chi2^{-\delta w}$.
Products of these factors are bounded by
$\prod_j(1+x_j)\le(1-\sum_jx_j)^{-1}$ when $x_j\ge0$ and
$\sum_jx_j<1$. The certified endpoint inequalities give a uniform mass
factor $M$. For a phase lower endpoint $a\le W$, define
\begin{align}
 M_1(a)&=\frac1{1-\chi(2^{-\delta q_0a}+E(\delta\theta_2a))},\notag\\
 M_2(a)&=\frac1{1-\chi E(\delta\theta_2a)},&
 M_3(a)&=\frac1{1-\chi E(\delta a)}.\label{eq:phase-masses}
\end{align}
All denominators are positive. In particular
\begin{equation}\label{eq:last-core-masses}
 S_{m_0-1}\le M_1(a)n,\qquad S_{m_0-2}\le M_2(a)n.
\end{equation}

\begin{lemma}[Stopped remainder]\label{lem:stopped}
For $a\le W\le d$, stop the actual partition before stage $k+1=m_0-2$.
Merge its frontier pieces on the original vertex set and call the resulting
graph $H$. Then
\[
     2\abs{E(H)}\le M_3(a)2^{\theta_2d}n.
\]
The long cycles retained before this stop number at most
\begin{equation}\label{eq:Jtrim}
                   J_{\rm tr}(a)n,
       \qquad J_{\rm tr}(a)=\tfrac12M_3(a)E(\ep a).
\end{equation}
\end{lemma}
\begin{proof}
The frontier average-degree bound is $2^{w_{k+1}}=2^{\theta_2W}$
per unit support mass. Its support mass is at most $M_3(a)n$.
The pieces have disjoint edges, so summing their degree bounds proves the
first inequality after the merge. Sum the long-cycle estimates from
Lemma~\ref{lem:scale}(i) over stages $0,\ldots,k$; their scales are bounded
below in reverse order by $a\theta_2^{-j}$. The same mass bound gives
\eqref{eq:Jtrim}.
\end{proof}

Every retained even core also has a completion packing of weight at most
$(1+C\eta)$ times its support size, provided the weighted order bound is
available for strictly smaller graphs. The completion construction uses
an edge bank, actual paths, and a smaller-order remainder. We state the
more precise version, including the extra credit needed for folding, next.

\section{Routing, completion, and original-edge ownership}\label{sec:routing}
This section records the graph-theoretic interfaces used by the physical
fold. They strengthen the routing and completion framework
of~\cite{OAI2026} quantitatively and keep its original-edge interpretation.
For a graph $H$, cut expansion at least $h$ means
\[
 e_H(Z,V(H)\setminus Z)\ge h\min\{\abs Z,\abs{V(H)\setminus Z}\}
                         \quad(Z\subseteq V(H)).
\]
For a supported core, cuts are taken inside its support.

\subsection{Two edge banks and smaller-order completion}
An even core at scale $w\ge w_0$ has expansion at least
$(1023/1024)2^{sw}$. A biased edge partition, with bias $511/512$ and
deviation $1/1024$, splits it into disjoint routing and completion banks
of expansions at least
\begin{equation}\label{eq:banks}
        \frac{1044483}{1048576}\,2^{sw},\qquad
        \frac{1023}{1048576}\,2^{sw},
\end{equation}
respectively. The cut-probability bound and its logarithmic scale gate
are certified uniformly for core orders at most $2^{\rho w}$.
There is no second parity correction: the cores were already corrected
in Lemma~\ref{lem:scale}.

For completion, the sampling-and-closing construction produces at most
one cycle per core vertex and leaves a graph supported on fewer than
$\eta$ times as many vertices. Apply the smaller-order induction hypothesis
to this induced residual graph and extend the packing by the constructed
cycles. Its weight is at most $(1+C\eta)$ times the core order. The strict
support inequality, rather than a comparison of average degrees, justifies
this induction call.

A fold also leaves high-degree external centres attached to a core. The
following is the precise credit needed in this case.
\begin{lemma}[High-centre completion]\label{lem:completion}
Let a graph be supported on $W=U\cup Z$, with $U\cap Z=\varnothing$ and
$2\le\abs W$. Suppose $w\ge w_A$, $2^{sw}\le\abs U\le2^{\rho w}$,
the graph has cut expansion on $U$ at least the second quantity in
\eqref{eq:banks}, and every external vertex has even degree and at least
$d$ neighbours in $U$. Suppose also
\[
 a_w=2^{\alpha w},\quad
 \deg(v,Z)\le a_w\ (v\in U),\quad
 \sqrt{2a_w\abs U}\le d\le\sqrt{2a_w\abs U}+1.
\]
If the weighted order bound with constant $K\ge0$ holds for all graphs of
strictly smaller order than $\abs W$, then the graph has a packing with
\begin{equation}\label{eq:high-credit}
 \wt_\sigma(\mathcal P)+d/2+2a_w
            \le(1+\sigma+K\eta+1/2048)\abs U.
\end{equation}
\end{lemma}
\begin{proof}
Put $q=\abs Z$ and $v=\abs U$. Double-counting edges from $Z$ to $U$
gives $qd\le a_wv$, and $d^2\ge2a_wv$ gives $q\le d/2$.
The scale inequalities give the stronger rounding budget
\[
                    q+d/2+2a_w\le v/2048.
\]
They also allow the core expansion to extend across the high centres with
the effective expansion required for completion. The completion construction
then gives weight at most
\[
                 (v+q)+(K\eta+\sigma)v.
\]
Only vertices of $U$ need a parity charge: vertices of $Z$ already have
even degree. Add $d/2+2a_w$ and use the rounding budget to obtain
\eqref{eq:high-credit}. The proof applies on the actual support $W$;
isolated padding is removed by a cost-preserving embedding. All residual
induction calls have fewer vertices than $W$.
\end{proof}

\subsection{A single routing system for all requests}
At an active scale put
\begin{equation}\label{eq:teams}
                 Q(w)=\ceil{c(w)2^{\alpha w}}.
\end{equation}
Requests consist of an ordered pair of endpoints, a team index, a set of
forbidden interior vertices, and a nonnegative vertex weight. Endpoint
load counts occurrences, so repeated endpoints are retained.

\begin{lemma}[Grouped weighted routing]\label{lem:routing}
In an even core at scale $w\ge w_A$, suppose that each vertex has endpoint
load at most $2^{\alpha w}$, each team contains at most $Q(w)$ requests,
and each forbidden set has size at most $2Q(w)$. There are paths in its
routing bank for all requests such that distinct paths are edge-disjoint;
paths in one team have disjoint interiors; each interior avoids its
forbidden set; and for every request with weight function $\omega$,
\begin{equation}\label{eq:soft-route}
       \sum_{v\text{ in its interior}}\omega(v)
                 \le\frac\gamma{Q(w)}\sum_{v\in U}\omega(v).
\end{equation}
Their lengths satisfy the scale-dependent bound
\mcode{globalRoutingLength} used in the capacity certificates.
\end{lemma}
\begin{proof}
For the large-ratio case, write $v$ for core order, $h$ for the routing
expansion, $B=Q(w)$ for team capacity and $k_0$ for endpoint capacity.
Put $q=1/c_0$, $a=a_{\rm prune}$, $R=7519120361/500000000$ and
$L=30589490361/5000000000$. The exact rational parameters satisfy
\begin{gather*}
 1+q+aq<2a^2(q+2),\qquad 4(1+q)/a<R,\\
 r_0:=1-a-1/128-1/(R\gamma)>0,\qquad
 8/r_0<(L-1/32)^2.
\end{gather*}
The scale inequalities give $k_0\le qB$ and
$R^2L^2B^2v\le h^3$. With $\ell=\floor{L\sqrt{v/h}}$ this implies
$RB\ell\le h$.

Associate to a candidate path its edge set and the marked interior vertices
$(\text{team},v)$. These are resources in
$E(G)\sqcup(\mathcal T\times V(G))$. A path of length at most $\ell$
uses at most $2\ell$ resources. To apply the Aharoni--Haxell selection
criterion~\cite{AH2000}, it suffices to show that for every nonempty
request subfamily $I$ and every resource set $Z$ of size less than
$4\ell\abs I$, some request in $I$ has an admissible path avoiding $Z$.

Write $F$ for the edge resources of $Z$ and $Y$ for its marked-vertex
resources. Put $D=4\ell\abs I$ and $x=\abs F/D$. Adaptive pruning chooses
$z,\tau>0$ with
\[
 z+\tau=a,\qquad \frac{qx}{z}+\frac{1-x}{\tau}=\frac{1+q}{a},
 \qquad \frac{qax}{2(1+q)}<\frac{z^2}{1+z}.
\]
After deletion of $F$, pruning discards at most $\abs F/(zh)$ vertices
and leaves expansion $(1-z)h$. Endpoint and team loads imply
\[
          \frac{k_0\abs F}{zh}+\frac{B\abs Y}{\tau h}<\abs I.
\]
Thus a request remains whose endpoints survive and whose team has at most
$\tau h$ marked interior obstructions. Delete those marks together with
its hard ports, and construct a short path in the surviving expanding core.
The short-path fan and averaging argument gives the additional weight bound
$(R\gamma\ell/h)\sum\omega\le(\gamma/B)\sum\omega$.
The retained expansion $r_0h$ gives the required path length $\ell$.
This proves the resource-avoidance condition. Aharoni--Haxell then selects
all paths simultaneously with disjoint resources, giving both forms of
disjointness in the statement. The formal selection lemma keeps repeated
endpoint occurrences as separate requests.

The actual integer $Q(w)$ is used throughout. Its rounding inequalities
prove the endpoint, team, hard-port and cubic-capacity gates at $w_A$ and
extend them to the full scale half-line.

There are two size regimes. If core order divided by routing-bank expansion
is at least $4096$, use the grouped weighted router directly. Otherwise
use the comparison bank and its bounded-ratio routing theorem. The
comparison ratio remains at least $4$ after the original parity correction.
Both regimes have the same conclusion~\eqref{eq:soft-route}. Thus no
intermediate range of core orders is omitted.

The completion machinery additionally uses an indexed reservoir router.
Its explicit gates are $z\ge8192$, $8192\ln v\le z$,
$2112\abs I\le z^2$, minimum degree at least $15z/16$, expansion at least
$z/30$, and endpoint load at most $z/64$. A balanced bank split leaves
one expanding bank to connect endpoint occurrences to a weak clique
immersion in the other. The clique-immersion theorem of~\cite{GLW2019}
and the terminal-pairing construction of~\cite{GMM2018}
then produce actual edge-disjoint paths. Only edge-disjointness is needed
from this reservoir theorem; it does not assert disjoint interiors.
All these graph constructions, including their integer capacities, are
proved in the formal source, not supplied by numerical gate checks alone.
\end{proof}

The key bookkeeping object for applying Lemma~\ref{lem:routing} is an
\emph{original endpoint certificate}. It assigns to each occurrence of
an endpoint of a request an incidence $(e,v)$ of an original edge, where
$v$ is that endpoint. The assignment is injective. If the original edges
used by two families are disjoint, their certificates combine. Reordering
the incoming endpoints within the same owner core preserves the certificate.

Three families are combined before routing: low two-edge paths, requests
from paired matching edges, and the switches of quotient cycles. Their
original edge sets are disjoint. The combined injection implies that the
\emph{total} endpoint load at $v$ is bounded by its degree in the assigned
descendant graph. Active-block preparation bounds this degree by
$2^{\alpha w}$. Bounding the three families separately and adding those
bounds would not establish the required load.

The combined team construction also preserves the $Q(w)$ capacity and
$2Q(w)$ hard-port limit. Apply Lemma~\ref{lem:routing} once to each selected
core, and use those same paths in all subsequent closures. Low groups and
paired matching groups close into actual simple cycles. Quotient cycles
are lifted through their switches, with the soft charge
\eqref{eq:soft-route}. The untouched completion bank remains available;
the core remainder deletes exactly the paths actually used. These ownership
rules are the reason that the four families in the next section have both
an exact edge cover and pairwise disjointness.

\section{Folding a layer and paying its cost}\label{sec:fold}
We work under the strong-induction hypothesis that every graph on fewer
than $n$ vertices has a packing of weight at most $C\phi(v)v$.
At an active stage $0\le i\le k$, put $w=w_i\ge w_A$ and
\[
 S=S_i,\qquad T=\sum_{j=0}^{i+2}S_j,\qquad f=\Phi(bw).
\]
Descendant pieces are grouped inside their actual DFS boxes into batches.
Write $D$ for total batch order and $A$ for the total size of the active
blocks that are contracted. The graph construction proves
\begin{equation}\label{eq:fold-masses}
 S-(N(w)+U(w))n\le A\le S+O(w)n,
       \qquad D\le(1+O(w))n.
\end{equation}
Here $N$ pays the small or unusable descendants, $U$ the discarded boundary
and grouping terms, and $O$ the overlapping support multiplicity.
Each batch has order at most $2^{bw}$. Its active blocks are pairwise
vertex-disjoint and belong to uniquely determined parent cores.

\begin{lemma}[Physical fold]\label{lem:physical}
There is an actual packing of the whole graph $G$ of weight at most
\begin{align}
 &(1+C\eta)T+\sigma S+S/2048+(B(w)+g(w))A\notag\\
 &\hspace{20mm}+Cf(D-\kappa A)+(J(w)+V(w))n.\label{eq:physical}
\end{align}
\end{lemma}
\begin{proof}
For each batch, retain one original representative of every quotient edge;
classify the original edges lost by contraction as internal pairs or
matching requests. This produces a simple quotient. An active block saves
at least $\kappa$ times its order in the quotient order, so the summed
quotient orders are at most $D-\kappa A$. Each quotient has order at most
$2^{bw}\le n^b<n$, since $n>1$ and $b<1$. Apply the induction hypothesis
and monotonicity of $\phi$ to obtain quotient packing weight at most
$Cf(D-\kappa A)$.

Use the joint routing construction of Section~\ref{sec:routing} on each
selected core. The single original-edge-endpoint injection covers low
pairs, paired matching edges, and switches in quotient cycles simultaneously.
The resulting four families are: completed core graphs, closed low groups,
closed matching groups, and lifted quotient packings. Uncovered quotient
edges keep their chosen original singleton representatives. Unused routing
paths are not deleted from the core.

The completed core and low-group families have combined weight at most
\[
 (1+\sigma+C\eta+1/2048)S
                  +\text{(prepared pair count)}/Q(w).
\]
This bound uses the high-centre credit of Lemma~\ref{lem:completion}
before any mass upper bound is substituted. The matching groups and the
prepared-pair fee together are at most $(B(w)+g(w))A$, and the soft routing
charge for lifted quotient cycles is at most $V(w)n$. The quotient packing
contributes the term already obtained above.

The four families are pairwise edge-disjoint and their union is exactly
the selected cores and the full descendant frontier. To see the importance
of both assertions, separate original non-routing edges from the paths
deleted during routing. Every original edge has exactly one of the four
owners by the quotient classification. Every deleted path has exactly its
request owner; joint routing makes different owners edge-disjoint. An edge
cannot belong to both classes because the core remainders were defined by
deleting precisely the selected paths. This is the partition proved in
\mcode{ActualFoldDisjoint12} and \mcode{ActualFoldCover12}.

Finally insert the other prefix cores, of total mass $T-S$, at cost
$(1+C\eta)(T-S)$, and the prefix long cycles at cost at most $J(w)n$.
The prefix partition from Section~\ref{sec:layers} now covers the whole
original graph. Adding the displayed costs proves~\eqref{eq:physical}.
\end{proof}

Define the coefficient of the selected core mass by
\begin{equation}\label{eq:fold-coefficient}
 \Gamma(P,w)=(1+C\eta)P+\sigma+B(w)+g(w)+\frac1{2048}
                                     -C\kappa\Phi(bw).
\end{equation}
\begin{lemma}[Fold ledger]\label{lem:ledger}
If $T\le PS$, then the right-hand side of~\eqref{eq:physical} is at most
\begin{equation}\label{eq:paid-fold}
                Cf n+\Gamma(P,w)S+CH(w)n.
\end{equation}
\end{lemma}
\begin{proof}
Write the terms involving $A$ as $BA-(C\kappa f-g)A$.
The second coefficient is nonnegative by~\eqref{eq:fee-endpoints}.
Consequently the \emph{lower} bound on $A$ in~\eqref{eq:fold-masses}
is used in its negative term, and the upper bound is used only in $BA$.
The resulting error coefficient of $n$ is
\[
 CfO+(C\kappa f-g)(N+U)+BO+J+V.
\]
All losses are nonnegative, $0<f\le1$, $g\ge0$, and $B\le B_0$, so this
is at most $CO+C\kappa(N+U)+B_0O+J+V=CH$.
The coefficient of $S$ is exactly~\eqref{eq:fold-coefficient}.
This proves~\eqref{eq:paid-fold}. In particular, one does not introduce
an extra $gO$ term by applying an upper bound to the negative group term.
\end{proof}

For each of the $40$ phase intervals $[a,d]$, the certificates supply
positive thresholds $P_t$ for all $t\ge3$. They prove
\begin{equation}\label{eq:fold-phase-payment}
 \Gamma(P_t,w)\le0\quad\text{whenever}\quad
                   w\ge a\theta_2^{-(t-3)}.
\end{equation}
There are explicit thresholds for $3\le t\le16$ and a constant threshold
for $t\ge16$. The scalar proof uses the monotonicity of $B,g,\Phi$;
the constant tail therefore requires only its first scale endpoint.

If $i$ is a stage with $S_i>0$ and
\begin{equation}\label{eq:fold-test}
                   \sum_{j=0}^{i+2}S_j\le P_{m_0-i}S_i,
\end{equation}
then $t=m_0-i$ and~\eqref{eq:schedule} give
$w_i=W\theta_2^{-(t-3)}\ge a\theta_2^{-(t-3)}$.
Lemmas~\ref{lem:physical} and~\ref{lem:ledger}, followed by
\eqref{eq:fold-phase-payment} and~\eqref{eq:potential-payment}, give
\[
        \wt_\sigma(\mathcal P)\le Cf n+CH(w_i)n\le C\phi(n)n.
\]
This completes every branch in which an active fold exists.

\section{The no-fold branch and the infinite tail}\label{sec:terminal}
Fix the phase interval $[a,d]$ containing $W$. Alongside $P_t$, its
certificate gives a positive summable sequence $(u_t)_{t\ge1}$ such that
\begin{equation}\label{eq:supersequence}
 u_1\ge M_1(a),\qquad u_2\ge M_2(a),\qquad
                  \sum_{s\ge t-2}u_s\le P_tu_t\quad(t\ge3).
\end{equation}
The initial terms are rational. The continuation is the positive recurrence
\begin{equation}\label{eq:tail-recurrence}
 u_{t+2}=(1-r)u_{t+1}+r(1-r)u_t\quad(t\ge15),
 \qquad r=\frac{88331750161}{10^{11}}.
\end{equation}
Since $0<r<1$, the recurrence is summable. If $\beta=r(1-r)$ and
$Q=1/[r^2(1-r)]$, then
\begin{equation}\label{eq:tail-sum}
       \sum_{s\ge15}u_s=Q\bigl(\beta u_{15}+(1-r)u_{16}\bigr).
\end{equation}
This follows by summing the recurrence. Its positive characteristic root
is less than $1$, or equivalently the same conclusion follows by the
geometric majorant used in the formal tail proof. The certificates check
the finite prefix of~\eqref{eq:supersequence} and the recurrence conditions
that extend it to every $t$; in particular $Q<P_{16}$.

\begin{lemma}[No-fold mass bound]\label{lem:no-fold-mass}
If no active stage with positive core mass satisfies~\eqref{eq:fold-test},
then
\[
                    S_i\le n u_{m_0-i}\quad(0\le i<m_0).
\]
In particular the cores retained before the terminal stop have total mass
at most $n\sum_{s=3}^{m_0}u_s$.
\end{lemma}
\begin{proof}
If the conclusion fails, take the maximum
$R=\max_i S_i/(nu_{m_0-i})>1$. The last two stages cannot attain this
maximum, by~\eqref{eq:last-core-masses} and~\eqref{eq:supersequence}.
At a maximizing active stage $i$, set $t=m_0-i\ge3$. Then
\[
 \sum_{j=0}^{i+2}S_j
   \le Rn\sum_{s=t-2}^{m_0}u_s
   \le RnP_tu_t=P_tS_i.
\]
Here $S_i>0$, so this is an active fold, a contradiction.
The retained stages are $0,\ldots,m_0-3$, giving the last assertion.
\end{proof}

Write
\[
 T_u=\sum_{s\ge3}u_s,\qquad R_m=\sum_{s\ge m+1}u_s.
\]
For the actual stopped graph $H$ in Lemma~\ref{lem:stopped}, the phase
certificate chooses one finite chain with fee $f_{a,d}$ and integer cap
at least $M_3(a)2^{\theta_2d}$. Thus $H$ has a packing of weight at most
$f_{a,d}n$. Packing the retained cores at cost $1+C\eta$ per unit mass
and adding the retained long cycles gives
\begin{equation}\label{eq:terminal-physical}
 \frac{\wt_\sigma(\mathcal P)}n
       \le f_{a,d}+(1+C\eta)(T_u-R_{m_0})+J_{\rm tr}(a).
\end{equation}
The negative tail $-R_{m_0}$ is kept throughout; it is an unused portion
of a summable majorant, not a negative physical graph cost.

\begin{lemma}[Terminal payment]\label{lem:terminal-payment}
For every phase and every integer $m\ge3$,
\begin{equation}\label{eq:terminal-payment}
 f_{a,d}+(1+C\eta)(T_u-R_m)+J_{\rm tr}(a)
                 \le C\Phi\bigl(a\theta_2^{-(m-3)}\bigr).
\end{equation}
\end{lemma}
\begin{proof}
The phase data first certify
\[
       f_{a,d}+(1+C\eta)T_u+J_{\rm tr}(a)\le C.
\]
For $3\le m<15$, they directly certify
\[
 C L\bigl(a\theta_2^{-(m-3)}\bigr)
    +f_{a,d}+(1+C\eta)T_u+J_{\rm tr}(a)
                      \le C+(1+C\eta)R_m.
\]
These are finite rational interval certificates for the positive series
in Section~\ref{sec:potential}.

For all later $m$, a stronger inequality suffices:
\begin{equation}\label{eq:zero-slack}
            C L\bigl(a\theta_2^{-(m-3)}\bigr)\le(1+C\eta)R_m.
\end{equation}
It is checked separately at $m=15$ and $m=16$. Positivity in
\eqref{eq:tail-recurrence} gives $R_{m+2}\ge\beta R_m$.
On the other hand Lemma~\ref{lem:loss-decay}, applied to every summand of
$L$, gives
\[
 L(x/\theta_2^2)\le
  2^{-(r_*/b)x(\theta_2^{-2}-1)}L(x).
\]
The certificate at $x=a\theta_2^{-12}$ bounds the exponential factor by
$\beta$. The factor only decreases as $x$ increases. Induction separately
on the two parities therefore proves~\eqref{eq:zero-slack} for every
$m\ge15$. Combine it with the total-budget inequality and use
$\Phi(x)=1-L(x)$ on this range to get~\eqref{eq:terminal-payment}.
\end{proof}

The $40$ closed phase intervals cover $[w_A,w_A/\theta_2]$. At a shared
endpoint either adjacent certificate can be used. For the actual phase,
$X=W\theta_2^{-(m_0-3)}\ge a\theta_2^{-(m_0-3)}$.
Equations~\eqref{eq:terminal-physical} and~\eqref{eq:terminal-payment},
and monotonicity of $\Phi$, show that the no-fold branch also has weight
at most $C\phi(n)n$.

\section{Completion of the proof}\label{sec:induction}
\begin{theorem}[Weighted potential bound]\label{thm:weighted-potential}
Every finite simple graph on $n$ vertices has a packing of weight at most
$C\phi(n)n$, where $C=485/42$ and $\phi$ is defined in~\eqref{eq:Phi}.
\end{theorem}
\begin{proof}
Use strong induction on $n$. If $2m\le2^{B_{\rm main}}n$, apply
Proposition~\ref{prop:finite-base}; this includes $n=0$.
Otherwise the simple-graph inequality $2m\le n^2$ implies
$n>2^{B_{\rm main}}$, and hence $X=\log_2n\ge w_A$.
Construct the actual scale decomposition of Section~\ref{sec:layers}.
All completion and quotient calls use graphs of strictly smaller order,
so the induction hypothesis supplies their packings.

If an active stage satisfies~\eqref{eq:fold-test}, apply the physical fold
and its scalar payment in Section~\ref{sec:fold}. Otherwise apply the
no-fold construction and terminal payment in Section~\ref{sec:terminal}.
In both cases the packing is on the original graph and has the claimed
weight. This completes the induction.
\end{proof}

\begin{proof}[Proof of Theorem~\ref{thm:main}]
Since $\phi(n)\le1$, Theorem~\ref{thm:weighted-potential} implies
\eqref{eq:weighted-order}. Proposition~\ref{prop:weighted-reduction}
then gives $\ce(G)\le12r(G)$, including the empty graph and isolated
components.
\end{proof}

\section{Formal statement and verification artifact}\label{sec:formal}
The formalization uses Lean~4~\cite{Lean4} and Mathlib~\cite{Mathlib}.
The declaration
\begin{center}\mcode{Cycle12.coefficient12}\end{center}
proves:
\begin{quote}\small\ttfamily
\textnormal{For every natural number \(n\) and simple graph \(G\) on
\(\mathrm{Fin}\,n\), there exists \(k\in\N\) such that}\\
Cycle13.EdgeDecomposition G k\\
\textnormal{and}\quad k $\le$ 12 * (n - Cycle13.componentCount G).
\end{quote}
The edge-decomposition predicate consists of a family of $k$ edge sets,
each a simple cycle edge set or a singleton original edge, with pairwise
disjointness and union equal to $G$'s edge set. The component count is the
cardinality of the type of all connected components. An additional inspection
example unfolds these predicates and writes the arithmetic with
\mcode{Nat.le}, \mcode{Nat.mul}, and \mcode{Nat.sub}. This checks the intended
statement, including isolated components, without any additional graph
hypothesis.

The supplied closure contains $10\,554$ project modules. The environment
is Lean $4.34.1$ and Mathlib revision
\begin{center}\small\mcode{d13f23b723b8a846827a245b89c10fc7d3f11612}.\end{center}
All project proof bodies in this release were recompiled from source in
this layout. After a temporary-storage reset, some lost object files were
recovered from historical caches only when their bytes matched the hashes
already recorded by this replacement build. Their prerequisite objects
were also required to match that recorded build; other objects were
regenerated. This recovery does not substitute a historical compilation
for a current-source compilation. The journal and recovery evidence
distinguish these operations. Preliminary replacement-interface builds
belong to the same run and use matching source, dependency and object hashes.
Fixed external Lean/Mathlib objects were reused; their sources were not
recompiled. Where import-header snapshots were used, an ordinary and a
snapshot-assisted compilation first produced byte-identical objects for
a current source. Every subsequent source body was still recompiled.
The final journal has a successful record for every retained module.

A traversal of the actual declaration types and proof bodies from the
final theorem found $\CycleDependencyCount$ dependencies, of which
$\CycleProjectDependencyCount$ belong to the retained project modules.
These include generated certificate lemmas, definitions, constructors and
recursors; they are not a count of separately authored mathematical ideas.
The retained modules contain no new axioms. The only axioms in the final
closure are
\[
              \mcode{propext},\quad\mcode{Classical.choice},\quad
              \mcode{Quot.sound}.
\]
No dependency is a \mcode{sorryAx}, unsafe declaration or partial declaration.


The historical declaration namespace \mcode{Cycle13} is retained for
compatibility. Current source paths are under \nolinkurl{Cycle12/Foundations}
and \nolinkurl{Cycle12/Independent}; this name does not change the constant
in \mcode{Cycle12.coefficient12}. Earlier compilation reports are historical
evidence, not substitutes for checking the replacement proofs.
The replacement graph bound follows Lemma~\ref{lem:eg-input}; the sampling
bound follows Lemma~\ref{lem:sampling-input}. Neither implementation imports
\mcode{ErdosProblems} or \mcode{ProbMethodCombinatorics}.
No empty-environment replay of the full declaration closure is claimed.

The full large dependency report is included as a gzip file, together with
its uncompressed SHA256 digest. The reproducibility scripts check every
source against the accepted manifest, rebuild the local closure in dependency
order, and check the exact theorem. They do not accept a file's existence
as evidence that it was compiled successfully. The repository documentation
separates reproducing the full proof from checking the packaged integrity
and inspecting an already built target.

The manuscript is a mathematical account of this formal proof, with its
principal correspondence listed in Appendix~\ref{app:correspondence}.
Lean verifies the formal declarations, not the correctness of this prose or
its bibliographic claims. The computational certificates in the source are
part of the proof artifact. Their generators and exploratory numerical
searches are not trusted proof oracles. This preprint does not claim
independent peer review or an optimal coefficient.

\appendix
\section{The short-circumference reduction and catalogue}\label{app:terminal}
We describe the computer-assisted proof of Lemma~\ref{lem:six} in enough
detail to specify both its infinite reduction and its finite verification.
For a nonempty connected graph the target is
\begin{equation}\label{eq:six-affine}
                    z(\mathcal P)\le19(n-1)+6m.
\end{equation}
A cut-vertex split is additive for this target: if its two sides meet only
at the cut vertex, their $(n_i-1)$ terms sum to $n-1$. Thus strong induction
reduces to a block that remains connected after deletion of any one vertex.
The small-order cases are checked directly. The corresponding result for
circumference at most $5$ is proved first by the same block induction.

If a minimal counterexample to~\eqref{eq:six-affine} has circumference $6$,
choose a six-cycle and label its vertices cyclically by $0,\ldots,5$.
The structural reduction uses ears between cycle vertices. An ear and the
longer of the two cycle arcs cannot form a cycle of length at least $7$.
The allowed ear endpoints and the compatibility of two ears are checked
with explicit path-splicing witnesses. Together with deletion of repeated
attachments, this shows that the graph outside the chosen cycle is a set
of isolated vertices together with at most one edge.

Here the induction reductions are essential. For example, two degree-two
vertices with the same pair of neighbours can be deleted and their four
edges restored as a four-cycle. Restoring costs $42$, whereas the target
allowance increases by $2\cdot19+4\cdot6=62$. The source also constructs
the analogous reductions for two exterior edge components with common
hubs and for three vertices with an alternating three-neighbour attachment.
They leave at most one vertex of each two-neighbour type and at most two
of each alternating three-neighbour type. These are reductions of actual
proper induced graphs, not assumptions that the whole graph has small order.

The $11$ possible attachment types, in their index order $0,\ldots,10$, are
\[
 \{0,2\},\{1,3\},\{2,4\},\{3,5\},\{0,4\},\{1,5\},
 \{0,3\},\{1,4\},\{2,5\},\{0,2,4\},\{1,3,5\}.
\]
These are the neighbourhoods in \mcode{SixAttachmentModels}. Pairwise incompatible types have explicit
simple long-cycle witnesses. All pairwise compatible type sets are
contained in one of the following $14$ families, using the type indices of
that file:
\[
\begin{gathered}
\{0,3,6\},\{1,3,6\},\{0,4,6\},\{1,4,6\},\quad
\{1,4,7\},\{2,4,7\},\{1,5,7\},\{2,5,7\},\\
\{0,2,8\},\{0,3,8\},\{2,5,8\},\{3,5,8\},\quad
\{0,2,4,9\},\{1,3,5,10\}.
\end{gathered}
\]
The first nine types have multiplicity at most $1$, and the last two at
most $2$. If an exterior edge is present, only the first four families
are possible after relabelling; its endpoints attach at opposite cycle
vertices $0,3$, with either or both of the two optional cross attachments.

The common model uses $16$ slots: six cycle vertices, four pairs of
attachment slots, and two slots for the optional exterior edge. At most
$11$ slots are active. The $12$ three-type families have $8$ multiplicity
choices each, and the last two families have $24$ choices each. There are
also $4\cdot8\cdot4$ nonempty exterior-edge choices. Hence the catalogue
has
\[
                       12\cdot8+2\cdot24+4\cdot8\cdot4=272
\]
parameter rows. The inverse index in \mcode{SixResidueParameterCoverage}
proves that every valid tuple is in this list. All nine non-cycle chords
of the six-cycle are optional, giving $272\cdot2^9=139\,264$ cases.

A certificate row is either a long-cycle word, or a list of pairwise
edge-disjoint cycle words. It records a set of required chords, and in
the packing case an upper bound $h$ on the total number of chords.
If $b_0$ is the number of base edges, $v_0$ the number of active vertices,
$c_0$ the number of certified cycles and $e_0$ their total number of
covered edges, the packing check is
\begin{equation}\label{eq:chord-check}
          42c_0+17(b_0+h)\le19(v_0-1)+23e_0.
\end{equation}
For any accepted chord set, $m'\le b_0+h$, so~\eqref{eq:chord-check}
implies
\[
                  42c_0+23(m'-e_0)\le19(v_0-1)+6m'.
\]
Every cycle word is checked for length, distinct vertices, consecutive
adjacencies and edge-disjointness. A long-cycle word excludes every chord
superset containing its required chords. The finite check proves that each
of the $512$ chord sets in each parameter row is accepted by one of its
certificates. All these checks are evaluated with \mcode{decide +kernel}.

Finally, actual outside vertices are enumerated in their attachment fibres.
The multiplicity bounds give an injection into the designated slots;
coverage gives a bijection to the active slots, and the classification gives
the complete graph equality, including the exterior edge in both
orientations. The finite witness transfers through this bijection.
The unused slots are isolated; removing them preserves the number of cycles,
the number of uncovered edges and the integer cost. This proves
\eqref{eq:six-affine}. Summing over connected components proves
Lemma~\ref{lem:six}, including isolated vertices and the graph of order zero.

\section{Certification of the mixed-cover construction}\label{app:cover}
This appendix specifies the hard case in Lemma~\ref{lem:mixed}. The
certificate concludes the existence of actual packings; it is not merely
a numerical upper bound for a hypothetical sample.

\subsection{Coordinates and fixed-size samples}
For $m>0$ put
\[
 x=\frac{\sqrt n}{B_*m^{1/3}},\quad s'=\frac{n^2}{2m}\ge1,\quad
 W(x)=B_*^3x^2,\quad v(x)=W(x)(1+3x).
\]
Then
\begin{equation}\label{eq:cover-coordinates}
 n=\frac{4B_*^3(s')^2}{W(x)^3},\qquad
 m=\frac{n^2}{2s'},\qquad D(n,m)/m=v(x).
\end{equation}
The irreducible range is $0<x\le17/250$, $s'\ge1$.
Set
\[
 p_0(x)=v(x)\left(\frac{99}{200}+\frac{11}{100}v(x)
                                      +\frac7{100}v(x)^2\right),\qquad
 N=\floor{p_0(x)n},\quad p=N/n.
\]
The sample is uniform among the $N$-element subsets of $V(G)$.
In particular, the proof does not substitute independent vertex sampling
for this fixed-cardinality law. For $P=p/W(x)$ and
$P_0(x)=p_0(x)/W(x)$ the exact floor bounds are
\[
 P_0(x)-\frac{W(x)^2}{4B_*^3(s')^2}\le P\le P_0(x).
\]
Monotone endpoint bounds turn these into rational bounds on every parameter
rectangle.

Let
\[
 H(x)=103/250+(351/200)x,\qquad D_0(x)=1-(33/125)W(x),
 \]\[
 h_G=B_*m^{1/3}H(x)/D_0(x),\qquad z=ph_G.
\]
Here $H(x)$ is a local coordinate function, unrelated to the loss $H(w)$
in Section~\ref{sec:potential}. The graph is assumed to meet the sparse-cut
irreducibility and minimum-degree inequalities left by the induction in
Lemma~\ref{lem:mixed}.

A good sample $U$ has size exactly $N$, certified lower neighbour degrees
into $U$, cut expansion inside $U$, an anchoring inequality, at least
$(97/100)p^2m$ induced edges, and at most $D(n,m)$ edges incident with $U$.
The anchoring inequality, for $X\subseteq V\setminus U$ and
$Y\subseteq U$ with $2\abs Y\le\abs U$, is
\[
 L\abs X\le e_G(X,U\setminus Y)+e_G(Y,U\setminus Y).
\]
Some rectangles also require a local cut bound up to sample proportion
$1/125$. The certified selection theorem bounds the union of six bad events:
low degree, a failed global cut, failed anchoring, too few sample edges,
too many incident edges, and the optional local-cut failure. Their rational
upper bounds have sum strictly less than $1$. Thus at least one actual
$N$-element sample is good. The fixed-size moment, degree, marked-cut and
anchor bounds and their parameter matches are proved before this union
bound is used.

\subsection{From the sample to cycles}
A capacity certificate splits the induced sample edges into banks, constructs
a weak clique immersion in the appropriate bank, and uses the remaining
expansion to route the requests from the path-and-cycle decomposition of Lov\'asz~\cite{Lovasz1968} applied to the
complement. Weak immersion is sufficient: routed paths are edge-disjoint;
branch vertices are not silently assumed to be avoided by all interiors.
The indexed routing and cycle-closure lemmas account for endpoint occurrences
and integer terminal capacities. The global-cut and additional-local-cut
branches have separate capacity inequalities.

The output is an actual packing $\mathcal C$ in $G$ with
\begin{equation}\label{eq:sample-output}
 p(\mathcal C)\le\frac{n-\abs U}{2},\qquad
 q(\mathcal C)\le D(n,m)+n-\abs U.
\end{equation}
The sample-size lower bound, including its floor, pays the fine cycle budget
$n\psi(m/(2B_*^3n))$. The second inequality in~\eqref{eq:sample-output}
pays both allowed residual-edge budgets, and the first is also at most
$n/2$. Thus this hard case proves both conclusions of
Lemma~\ref{lem:mixed}. The full induction may still choose different
packings for its two conclusions.

\subsection{Continuous coverage and its soundness interface}
A rational rectangle is a tuple $(x_-,x_+,s_-,s_+)$, where $s_+$ may be
absent, denoting an unbounded ray. Its \emph{soundness} statement quantifies
over every actual graph whose coordinates lie in it and which satisfies
the irreducibility and degree conditions, and concludes
\eqref{eq:fine-cover} and~\eqref{eq:coarse-cover}.
A routed case consists of a sample certificate and a capacity certificate.
Their common parameters---degree, expansion, anchoring and any optional
local cut constant---must agree exactly. The theorem
\mcode{rat_routed_case_sound} combines their graph-theoretic conclusions
only after these equalities and domain inclusions are checked.

Rectangles are joined by exact rational interval chains. Horizontal joining
checks that the $x$ intervals cover and that each piece extends over the
required $s'$ interval. Vertical joining checks the converse. Both operations
include closed endpoints, and an unbounded target ray must be covered by
a chain that actually reaches an unbounded piece. These are proved
soundness-preserving operations, rather than a sampled-grid argument.
The final assembly is
\[
 [0,3/1000]\times[1,\infty)\quad\cup\quad
 [3/1000,17/250]\times[1,\infty),
\]
with each rectangle further split into its certified scale pieces.
Its final statement is
\begin{center}\mcode{Cycle12.CoverAssembly.Full.sound}.\end{center} Actual graphs have $x>0$;
including $x=0$ in the interval certificate causes no division by zero in
the graph application.

This closes the irreducible case and hence the strong induction. The
certificate hierarchy contains the actual probability bounds and routing
capacities, not just the final interval endpoints. Its rational decisions
are checked in Lean, and its probabilistic and graph-theoretic soundness
lemmas lie in the same axiom-audited dependency closure as the main theorem.

\section{Exact parameters and certificate indices}\label{app:parameters}
The following values are read from \mcode{GlobalParameters12} and checked
for exact equality with the documentary input \mcode{paper/data/global12.json}.
Fractions in this table are exact. Large finite certificate entries remain
in the Lean source; the supplementary JSON files provide their indices.

\begin{longtable}{@{}ll@{}}
\toprule Parameter & Exact value\\\midrule\endhead
$w_0$ & $\frac{10170985991133}{100000000000}$\\[3pt]
$s$ & $\frac{38082365177}{50000000000}$\\[3pt]
$\theta_0$ & $\frac{2048519349}{2500000000}$\\[3pt]
$\theta_2$ & $\frac{15961744051}{20000000000}$\\[3pt]
$\alpha$ & $\frac{13908880001}{25000000000}$\\[3pt]
$\lambda$ & $\frac{3511913879}{10000000000}$\\[3pt]
$b$ & $\frac{3004140221}{4000000000}$\\[3pt]
$\gamma$ & $\frac{2460109537}{1562500000}$\\[3pt]
$\delta$ & $\frac{3454748071}{50000000000}$\\[3pt]
$\varepsilon$ & $\frac{3093377987}{100000000000}$\\[3pt]
$c_0$ & $\frac{284161497}{500000000}$\\[3pt]
$a_{\rm prune}$ & $\frac{36700631257}{50000000000}$\\[3pt]
$\theta_1$ & $\frac{10059807341}{12500000000}$\\[3pt]
$\rho$ & $\frac{110002874129}{100000000000}$\\[3pt]
$w_A$ & $\frac{2311587725257500000000}{14597460907219352581}$\\[3pt]
$\eta$ & $\frac{244357}{2000000000}$\\[3pt]
$M$ & $\frac{10218421973}{10000000000}$\\[3pt]
$\chi$ & $\frac{10077234539}{5000000000}$\\[3pt]
$B_0$ & $\frac{1225527173}{10000000000}$\\[3pt]
\mcode{router_length} & $\frac{30589490361}{5000000000}$\\[3pt]
\mcode{route_load} & $\frac{7519120361}{500000000}$\\[3pt]
\mcode{order_bound} & $\frac{7065355693}{10000000000}$\\[3pt]
$C$ & $\frac{485}{42}$\\[3pt]
$\sigma$ & $\frac{23}{42}$\\[3pt]
$\kappa$ & $\frac{999999}{1000000}$\\[3pt]
\mcode{edge_bias} & $\frac{511}{512}$\\[3pt]
\mcode{split_deviation} & $\frac{1}{1024}$\\[3pt]
$p$ & $\frac{244357}{8000000000}$\\[3pt]
\mcode{router_bank} & $\frac{1044483}{1048576}$\\[3pt]
\mcode{completion_bank} & $\frac{1023}{1048576}$\\[3pt]
\mcode{router_constant} & $\frac{52902787547485085406377454545012954383041}{6250000000000000000000000000000000000}$\\[3pt]
$\phi_0$ & $\frac{46220729}{50000000}$\\[3pt]
\bottomrule\end{longtable}
The archived ordinary-proof finite inputs are included for provenance. Their
reported upper bounds can differ in the last few decimal places from the
formally proved chain constants. This paper and its tables use the Lean
constants exclusively; the index records the exact differences. No historical
Python fee is used as a premise of the formal proof.

The main degree logarithm is $158387/1000$. The tail parameter is
$r=88331750161/10^{11}$, with $\beta=r(1-r)$ and $Q=1/[r^2(1-r)]$.
The parameters \mcode{p}, \mcode{router_length}, \mcode{route_load},
\mcode{order_bound}, and \mcode{router_constant} belong to the completion
and routing gates; they are not independent choices left open in the proof.

Rational certificates enclose powers of $2$, logarithms, square roots and
cube roots. For a positive series, a finite partial sum is a lower bound;
an upper bound includes a proved geometric tail. The renewal loss is
bounded by a positive finite expansion with its proved tail, and all
inequalities are preserved in the required direction. Tangent certificates
control continuous density intervals; monotonicity and rate certificates
control scale half-lines. These steps justify extending the finite data to
all graphs and all numbers of stages.

The phase table below is for orientation: its displayed phase endpoints
are rounded to six decimal places and are \emph{not} used as certificate
endpoints. The exact endpoints, thresholds $P_3,\ldots,P_{16}$, and all
$u_1,\ldots,u_{16}$ are in \mcode{paper/data/phase-base-index.json} and in
the corresponding Lean modules. The listed fees are the actual rational
\mcode{costBound} constants of the Lean finite chains.

\begin{longtable}{@{}rrrrr@{}}
\toprule Phase & $a$ (display) & $d$ (display) & Chain & $f_{a,d}$\\\midrule\endhead
00 & 158.355466 & 159.607444 & 011 & $\frac{5391841825851}{500000000000}$\\[3pt]
01 & 159.607444 & 160.859421 & 069 & $\frac{10796623618537}{1000000000000}$\\[3pt]
02 & 160.859421 & 162.111399 & 012 & $\frac{2702389777377}{250000000000}$\\[3pt]
03 & 162.111399 & 163.363377 & 070 & $\frac{2705625297969}{250000000000}$\\[3pt]
04 & 163.363377 & 164.615354 & 014 & $\frac{2708860390719}{250000000000}$\\[3pt]
05 & 164.615354 & 165.867332 & 071 & $\frac{10848385314567}{1000000000000}$\\[3pt]
06 & 165.867332 & 167.119310 & 015 & $\frac{10861328073341}{1000000000000}$\\[3pt]
07 & 167.119310 & 168.371287 & 017 & $\frac{2718571646969}{250000000000}$\\[3pt]
08 & 168.371287 & 169.623265 & 043 & $\frac{10887234417859}{1000000000000}$\\[3pt]
09 & 169.623265 & 170.875243 & 000 & $\frac{10900183002579}{1000000000000}$\\[3pt]
10 & 170.875243 & 172.127220 & 073 & $\frac{1364141621327}{125000000000}$\\[3pt]
11 & 172.127220 & 173.379198 & 019 & $\frac{2731521075901}{250000000000}$\\[3pt]
12 & 173.379198 & 174.631176 & 059 & $\frac{1367379994851}{125000000000}$\\[3pt]
13 & 174.631176 & 175.883153 & 021 & $\frac{10951994641041}{1000000000000}$\\[3pt]
14 & 175.883153 & 177.135131 & 022 & $\frac{10959542293243}{1000000000000}$\\[3pt]
15 & 177.135131 & 178.387109 & 023 & $\frac{10966965799873}{1000000000000}$\\[3pt]
16 & 178.387109 & 179.639086 & 024 & $\frac{2743615751629}{250000000000}$\\[3pt]
17 & 179.639086 & 180.891064 & 044 & $\frac{10982037016331}{1000000000000}$\\[3pt]
18 & 180.891064 & 182.143042 & 026 & $\frac{10989687427511}{1000000000000}$\\[3pt]
19 & 182.143042 & 183.395019 & 076 & $\frac{10997418068539}{1000000000000}$\\[3pt]
20 & 183.395019 & 184.646997 & 051 & $\frac{11005220797637}{1000000000000}$\\[3pt]
21 & 184.646997 & 185.898975 & 052 & $\frac{1376639403931}{125000000000}$\\[3pt]
22 & 185.898975 & 187.150952 & 061 & $\frac{1377635136691}{125000000000}$\\[3pt]
23 & 187.150952 & 188.402930 & 032 & $\frac{11029125826993}{1000000000000}$\\[3pt]
24 & 188.402930 & 189.654908 & 077 & $\frac{2759313584727}{250000000000}$\\[3pt]
25 & 189.654908 & 190.906885 & 034 & $\frac{11045464831563}{1000000000000}$\\[3pt]
26 & 190.906885 & 192.158863 & 063 & $\frac{138172005771}{12500000000}$\\[3pt]
27 & 192.158863 & 193.410841 & 064 & $\frac{11062149004901}{1000000000000}$\\[3pt]
28 & 193.410841 & 194.662818 & 035 & $\frac{11070615823097}{1000000000000}$\\[3pt]
29 & 194.662818 & 195.914796 & 079 & $\frac{2769792387169}{250000000000}$\\[3pt]
30 & 195.914796 & 197.166774 & 038 & $\frac{277194651861}{25000000000}$\\[3pt]
31 & 197.166774 & 197.792762 & 039 & $\frac{5546050006983}{500000000000}$\\[3pt]
32 & 197.792762 & 198.105757 & 081 & $\frac{11094248594251}{1000000000000}$\\[3pt]
33 & 198.105757 & 198.262254 & 040 & $\frac{11095327100827}{1000000000000}$\\[3pt]
34 & 198.262254 & 198.340503 & 082 & $\frac{1109587104771}{100000000000}$\\[3pt]
35 & 198.340503 & 198.379627 & 067 & $\frac{5548069169719}{500000000000}$\\[3pt]
36 & 198.379627 & 198.399189 & 041 & $\frac{11096276645863}{1000000000000}$\\[3pt]
37 & 198.399189 & 198.408970 & 065 & $\frac{2219267420263}{200000000000}$\\[3pt]
38 & 198.408970 & 198.413861 & 047 & $\frac{554818554329}{50000000000}$\\[3pt]
39 & 198.413861 & 198.418751 & 066 & $\frac{2774101376139}{250000000000}$\\[3pt]
\bottomrule\end{longtable}
The $45$ base intervals and their assigned chains are listed in the same
exact index file. The following table records all $83$ finite chains. Its
cutoff $x$ means the average-degree condition $2m\le2^xn$; the file also
records the larger integer degree cap used for this condition. The column
$u$ is the denominator of the harmonic starting bound, and $s$ is the
number of certified DFS steps. Every chain uses the same mixed-cover
parameters $B_*=43/10$, $A_*=3$ and singleton weight $23/42$.

\begin{longtable}{@{}rrrrr@{}}
\toprule Chain & Exact cutoff $x$ & $s$ & $u$ & Certified fee\\\midrule\endhead
000 & $\frac{1091}{8}$ & 0 & 224502 & $\frac{10900183002579}{1000000000000}$\\[3pt]
001 & $\frac{1129}{8}$ & 0 & 117478 & $\frac{2191544772857}{200000000000}$\\[3pt]
002 & $\frac{118931}{1000}$ & 0 & 144130 & $\frac{10674267180743}{1000000000000}$\\[3pt]
003 & $\frac{118933}{1000}$ & 0 & 144173 & $\frac{2134858533271}{200000000000}$\\[3pt]
004 & $\frac{118941}{1000}$ & 0 & 144240 & $\frac{10674395954517}{1000000000000}$\\[3pt]
005 & $\frac{11897}{100}$ & 0 & 144304 & $\frac{2134954301487}{200000000000}$\\[3pt]
006 & $\frac{119239}{1000}$ & 0 & 145330 & $\frac{5339126553537}{500000000000}$\\[3pt]
007 & $\frac{121397}{1000}$ & 0 & 153684 & $\frac{10706185307421}{1000000000000}$\\[3pt]
008 & $\frac{12263}{100}$ & 0 & 157998 & $\frac{2680536955283}{250000000000}$\\[3pt]
009 & $\frac{123863}{1000}$ & 0 & 163280 & $\frac{2684527045207}{250000000000}$\\[3pt]
010 & $\frac{126329}{1000}$ & 0 & 173733 & $\frac{1346254529237}{125000000000}$\\[3pt]
011 & $\frac{127383}{1000}$ & 0 & 178559 & $\frac{5391841825851}{500000000000}$\\[3pt]
012 & $\frac{129381}{1000}$ & 0 & 187995 & $\frac{2702389777377}{250000000000}$\\[3pt]
013 & $\frac{131261}{1000}$ & 0 & 197117 & $\frac{10833912573677}{1000000000000}$\\[3pt]
014 & $\frac{131379}{1000}$ & 0 & 197460 & $\frac{2708860390719}{250000000000}$\\[3pt]
015 & $\frac{133377}{1000}$ & 0 & 208051 & $\frac{10861328073341}{1000000000000}$\\[3pt]
016 & $\frac{133727}{1000}$ & 0 & 209257 & $\frac{2173173368277}{200000000000}$\\[3pt]
017 & $\frac{134377}{1000}$ & 0 & 213451 & $\frac{2718571646969}{250000000000}$\\[3pt]
018 & $\frac{136193}{1000}$ & 0 & 223406 & $\frac{5448912124509}{500000000000}$\\[3pt]
019 & $\frac{138373}{1000}$ & 0 & 236415 & $\frac{2731521075901}{250000000000}$\\[3pt]
020 & $\frac{138659}{1000}$ & 0 & 237842 & $\frac{10929793693067}{1000000000000}$\\[3pt]
021 & $\frac{140371}{1000}$ & 0 & 248776 & $\frac{10951994641041}{1000000000000}$\\[3pt]
022 & $\frac{141371}{1000}$ & 0 & 117483 & $\frac{10959542293243}{1000000000000}$\\[3pt]
023 & $\frac{14237}{100}$ & 0 & 117477 & $\frac{10966965799873}{1000000000000}$\\[3pt]
024 & $\frac{143369}{1000}$ & 0 & 117479 & $\frac{2743615751629}{250000000000}$\\[3pt]
025 & $\frac{143591}{1000}$ & 0 & 117477 & $\frac{1097613966829}{100000000000}$\\[3pt]
026 & $\frac{145367}{1000}$ & 0 & 117459 & $\frac{10989687427511}{1000000000000}$\\[3pt]
027 & $\frac{146057}{1000}$ & 0 & 117478 & $\frac{171797133157}{15625000000}$\\[3pt]
028 & $\frac{14729}{100}$ & 0 & 117476 & $\frac{1375579270293}{125000000000}$\\[3pt]
029 & $\frac{148523}{1000}$ & 0 & 117477 & $\frac{11014368472807}{1000000000000}$\\[3pt]
030 & $\frac{14867}{125}$ & 0 & 144185 & $\frac{10674331494659}{1000000000000}$\\[3pt]
031 & $\frac{14876}{125}$ & 0 & 144442 & $\frac{5337631723901}{500000000000}$\\[3pt]
032 & $\frac{150363}{1000}$ & 0 & 117493 & $\frac{11029125826993}{1000000000000}$\\[3pt]
033 & $\frac{150989}{1000}$ & 0 & 117474 & $\frac{11034209376807}{1000000000000}$\\[3pt]
034 & $\frac{152361}{1000}$ & 0 & 117445 & $\frac{11045464831563}{1000000000000}$\\[3pt]
035 & $\frac{155359}{1000}$ & 0 & 117478 & $\frac{11070615823097}{1000000000000}$\\[3pt]
036 & $\frac{155921}{1000}$ & 0 & 117479 & $\frac{2768854402683}{250000000000}$\\[3pt]
037 & $\frac{15637}{125}$ & 0 & 168384 & $\frac{10754072517109}{1000000000000}$\\[3pt]
038 & $\frac{157357}{1000}$ & 0 & 118851 & $\frac{277194651861}{25000000000}$\\[3pt]
039 & $\frac{157857}{1000}$ & 0 & 119873 & $\frac{5546050006983}{500000000000}$\\[3pt]
040 & $\frac{158231}{1000}$ & 0 & 120651 & $\frac{11095327100827}{1000000000000}$\\[3pt]
041 & $\frac{158341}{1000}$ & 0 & 120834 & $\frac{11096276645863}{1000000000000}$\\[3pt]
042 & $\frac{158387}{1000}$ & 0 & 120962 & $\frac{11096672924573}{1000000000000}$\\[3pt]
043 & $\frac{16922}{125}$ & 0 & 218645 & $\frac{10887234417859}{1000000000000}$\\[3pt]
044 & $\frac{18046}{125}$ & 0 & 117477 & $\frac{10982037016331}{1000000000000}$\\[3pt]
045 & $\frac{18103}{125}$ & 0 & 117478 & $\frac{2746379951309}{250000000000}$\\[3pt]
046 & $\frac{19336}{125}$ & 0 & 117477 & $\frac{5532459781859}{500000000000}$\\[3pt]
047 & $\frac{19794}{125}$ & 0 & 120882 & $\frac{554818554329}{50000000000}$\\[3pt]
048 & $\frac{2379}{20}$ & 0 & 144231 & $\frac{10674512940043}{1000000000000}$\\[3pt]
049 & $\frac{23817}{200}$ & 0 & 144740 & $\frac{1067625989843}{100000000000}$\\[3pt]
050 & $\frac{25759}{200}$ & 0 & 184710 & $\frac{1350246585773}{125000000000}$\\[3pt]
051 & $\frac{29473}{200}$ & 0 & 117438 & $\frac{11005220797637}{1000000000000}$\\[3pt]
052 & $\frac{29673}{200}$ & 0 & 117478 & $\frac{1376639403931}{125000000000}$\\[3pt]
053 & $\frac{29733}{250}$ & 0 & 144141 & $\frac{1067427985913}{100000000000}$\\[3pt]
054 & $\frac{29887}{250}$ & 0 & 146490 & $\frac{10682252366379}{1000000000000}$\\[3pt]
055 & $\frac{30041}{250}$ & 0 & 148859 & $\frac{5345112523633}{500000000000}$\\[3pt]
056 & $\frac{30691}{200}$ & 0 & 117472 & $\frac{1105455351259}{100000000000}$\\[3pt]
057 & $\frac{32507}{250}$ & 0 & 190826 & $\frac{10817943414113}{1000000000000}$\\[3pt]
058 & $\frac{3374}{25}$ & 0 & 216265 & $\frac{5440921458393}{500000000000}$\\[3pt]
059 & $\frac{34843}{250}$ & 0 & 242342 & $\frac{1367379994851}{125000000000}$\\[3pt]
060 & $\frac{34973}{250}$ & 0 & 245655 & $\frac{437831386747}{40000000000}$\\[3pt]
061 & $\frac{37341}{250}$ & 0 & 117473 & $\frac{1377635136691}{125000000000}$\\[3pt]
062 & $\frac{37439}{250}$ & 0 & 117480 & $\frac{11024226683983}{1000000000000}$\\[3pt]
063 & $\frac{3834}{25}$ & 0 & 117464 & $\frac{138172005771}{12500000000}$\\[3pt]
064 & $\frac{3859}{25}$ & 0 & 117477 & $\frac{11062149004901}{1000000000000}$\\[3pt]
065 & $\frac{39587}{250}$ & 0 & 120849 & $\frac{2219267420263}{200000000000}$\\[3pt]
066 & $\frac{39589}{250}$ & 0 & 120885 & $\frac{2774101376139}{250000000000}$\\[3pt]
067 & $\frac{6333}{40}$ & 0 & 120808 & $\frac{5548069169719}{500000000000}$\\[3pt]
068 & $\frac{63781}{500}$ & 0 & 179618 & $\frac{10786001135279}{1000000000000}$\\[3pt]
069 & $\frac{64191}{500}$ & 0 & 182773 & $\frac{10796623618537}{1000000000000}$\\[3pt]
070 & $\frac{6519}{50}$ & 0 & 192519 & $\frac{2705625297969}{250000000000}$\\[3pt]
071 & $\frac{66189}{500}$ & 0 & 202782 & $\frac{10848385314567}{1000000000000}$\\[3pt]
072 & $\frac{66247}{500}$ & 0 & 203277 & $\frac{5424943263449}{500000000000}$\\[3pt]
073 & $\frac{68687}{500}$ & 0 & 230378 & $\frac{1364141621327}{125000000000}$\\[3pt]
074 & $\frac{68713}{500}$ & 0 & 230694 & $\frac{109138064341}{10000000000}$\\[3pt]
075 & $\frac{71179}{500}$ & 0 & 117477 & $\frac{10966874863597}{1000000000000}$\\[3pt]
076 & $\frac{73183}{500}$ & 0 & 117475 & $\frac{10997418068539}{1000000000000}$\\[3pt]
077 & $\frac{75681}{500}$ & 0 & 117480 & $\frac{2759313584727}{250000000000}$\\[3pt]
078 & $\frac{76111}{500}$ & 0 & 117476 & $\frac{2761079337139}{250000000000}$\\[3pt]
079 & $\frac{78179}{500}$ & 0 & 117507 & $\frac{2769792387169}{250000000000}$\\[3pt]
080 & $\frac{78577}{500}$ & 0 & 118427 & $\frac{11086034684193}{1000000000000}$\\[3pt]
081 & $\frac{79053}{500}$ & 0 & 120382 & $\frac{11094248594251}{1000000000000}$\\[3pt]
082 & $\frac{79147}{500}$ & 0 & 120696 & $\frac{1109587104771}{100000000000}$\\[3pt]
\bottomrule\end{longtable}

\section{Correspondence with Lean}\label{app:correspondence}
The following declarations are present in the audited transitive proof-body
closure of the main theorem. The machine-readable table
\mcode{docs/proof-map.json} records their modules and source paths. The
conditional induction interface in the last part of the proof is instantiated
with the proved physical construction; it is not an additional assumption
of Theorem~\ref{thm:main}. Some source comments record earlier intermediate
obligations. Their discharge is determined by the final theorem and its
actual dependency chain, not by those historical comments.

\begin{longtable}{@{}p{0.30\textwidth}p{0.66\textwidth}@{}}
\toprule Statement & Lean declaration\\\midrule\endhead
Final original-edge partition (\ref{thm:main}) & {\small\nolinkurl{Cycle12.coefficient12}}\\[7pt]
Rotation and deletion cycle bound (\ref{lem:eg-input}) & {\small\nolinkurl{Cycle12.Independent.erdosGallai_cycle}}\\[7pt]
Bernoulli sample lower tail (\ref{lem:sampling-input}) & {\small\nolinkurl{Cycle13.vertex_sample_edge_lower_tail}}\\[7pt]
Weighted order to coefficient twelve (\ref{prop:weighted-reduction}) & {\small\nolinkurl{Cycle12.coefficient12_from_weighted_order_bound}}\\[7pt]
Finite rooted amplification (\ref{lem:amplification}) & {\small\nolinkurl{Cycle12.rooted_weighted_order_to_predecessor}}\\[7pt]
Circumference at most six (\ref{lem:six}) & {\small\nolinkurl{Cycle12.terminal_rank_no_seven_cycle}}\\[7pt]
Harmonic affine extension (\ref{lem:harmonic}) & {\small\nolinkurl{Cycle13.weightedAffine_harmonic}}\\[7pt]
Mixed cover (\ref{lem:mixed}) & {\small\nolinkurl{Cycle12.mixed_cover}}\\[7pt]
Certified DFS step (\ref{lem:dfs-step}) & {\small\nolinkurl{Cycle12.FiniteStageCertificate.step}}\\[7pt]
Finite-degree potential bound (\ref{prop:finite-base}) & {\small\nolinkurl{Cycle12.Global12.finite_degree_potential}}\\[7pt]
Uniform loss decay (\ref{lem:loss-decay}) & {\small\nolinkurl{Cycle12.Global12.lossH_decay}}\\[7pt]
Renewal potential payment (\ref{lem:potential}) & {\small\nolinkurl{Cycle12.Global12.globalPhi_fold_budget}}\\[7pt]
Actual even scale layers (\ref{sec:layers}) & {\small\nolinkurl{Cycle12.Even12.exists_actual_scale_layers}}\\[7pt]
High-centre completion (\ref{lem:completion}) & {\small\nolinkurl{Cycle12.Even12.high_center_completion_at_scale}}\\[7pt]
Even-bank grouped weighted routing (\ref{lem:routing}) & {\small\nolinkurl{Cycle12.Even12.bank_routing_at_scale}}\\[7pt]
Joint original endpoint family (\ref{sec:routing}) & {\small\nolinkurl{Cycle12.Even12.ActualFoldStage12.jointFamily}}\\[7pt]
Active mass lower bound (\ref{eq:fold-masses}) & {\small\nolinkurl{Cycle12.Even12.ActualFoldStage12.activeMass_lower}}\\[7pt]
Four-family actual packing payment (\ref{lem:physical}) & {\small\nolinkurl{Cycle12.Even12.ActualFoldStage12.actual_fold_target_payment}}\\[7pt]
Whole-graph physical construction (\ref{lem:physical}) & {\small\nolinkurl{Cycle12.Global12.physical_fold_construction12}}\\[7pt]
Physical scalar ledger (\ref{lem:ledger}) & {\small\nolinkurl{Cycle12.Global12.physical_fold_ledger_payment}}\\[7pt]
No-fold core mass bound (\ref{lem:no-fold-mass}) & {\small\nolinkurl{Cycle12.Even12.EvenScaleLayers12.kept_mass_of_no_fold}}\\[7pt]
Infinite terminal tail (\ref{lem:terminal-payment}) & {\small\nolinkurl{Cycle12.Global12.infinite_terminal_tail_payment}}\\[7pt]
Terminal potential payment (\ref{lem:terminal-payment}) & {\small\nolinkurl{Cycle12.Global12.terminal_potential_payment}}\\[7pt]
Strong induction from physical branch (\ref{thm:weighted-potential}) & {\small\nolinkurl{Cycle12.Global12.potential_bound_from_fold_branch}}\\[7pt]
Coverage of finite parameter tuples (\ref{app:terminal}) & {\small\nolinkurl{Cycle12.six_residue_parameter_coverage}}\\[7pt]
Full continuous cover (\ref{app:cover}) & {\small\nolinkurl{Cycle12.CoverAssembly.Full.sound}}\\[7pt]
\bottomrule\end{longtable}

\bibliographystyle{amsplain}
\bibliography{references}
\section*{Acknowledgements and use of AI tools}
OpenAI ChatGPT/Codex tools were used extensively in the development of the
argument, computational exploration, preparation of the Lean formalization,
and drafting and organization of this manuscript. The source repository
records the formal theorem, its proof dependencies, and the scope of the
reported checks. These tools are not authors; responsibility for the submitted
work rests with the author. The author also acknowledges the contributors to
Lean, Mathlib, and the third-party formal libraries identified in the repository.
\end{document}
